\providecommand{\HH}{\operatorname{HH}}
\providecommand{\Har}{\operatorname{Har}}

\section{Bulk–boundary functoriality and spectral $R(z)$}
\label{sec:bulk-boundary-R}
\label{sec:spectral_braiding}
\label{sec:spectral-braiding}
\label{chap:spectral-braiding}
\label{sec:bulk-boundary}

The spectral $R$-matrix $R(z) = \Res^{\mathrm{coll}}_{0,2}(\Theta^{\mathrm{oc}})$ is the genus-0, arity-2 collision residue of the open/closed MC element, extracted from the bulk-to-boundary composition of two parallel line operators.  It satisfies the Yang--Baxter identity (Stokes on $\FM_3(\C)$) and equals the Laplace transform of the chiral $\lambda$-bracket.

\subsection{Boundary data and factorization module structure}
\label{subsec:boundary-module}
Let $A=(A_{\mathsf{ch}},A_{\mathsf{top}})$ be a $C_\ast\!\bigl(W(\mathsf{SC}^{\mathrm{ch,top}})\bigr)$–algebra.
A (topological) boundary condition along $t=0$ is modeled by a \emph{right $W(\mathsf{SC}^{\mathrm{ch,top}})$–module} $M$ supported on $\C\times \R_{\ge 0}$, i.e.\ a prefactorization algebra $\mathsf{Obs}^{\partial}$ on half‑rectangles with actions
\[
C_\ast\bigl(\FM_k(\C)\times E_1(m)\bigr)\otimes A_{\mathsf{ch}}^{\otimes k}\otimes M^{\otimes m}\longrightarrow M,
\]
compatible with $W$–coherence.  We refer to $M$ as a \emph{boundary factorization module}.

\subsection{Line operators and the spectral parameter}
\label{subsec:line-ops}
Boundary line operators $L$ are point‑supported defects along $t=0$ with insertion points $z\in\C$.
Given two parallel boundary lines $L_1(z_1)$ and $L_2(z_2)$ (with $z_1\neq z_2$), the bulk-to-boundary OPE in the strip determines a composition
\[
\mathsf{Comp}_{z_1,z_2} \colon L_1(z_1)\otimes L_2(z_2) \longrightarrow L_2(z_2)\otimes L_1(z_1),
\]
meromorphic in $z_{12}=z_1-z_2$ with logarithmic poles along $z_{12}=0$.
This composition is induced by the $W(\mathsf{SC}^{\mathrm{ch,top}})$–module structure and is computed by configuration integrals over $\FM_2(\C)$ with $E_1$–homotopies along the half‑interval.

\begin{definition}[Spectral braiding]
\label{def:spectral-braiding}
The \emph{spectral braiding} (or \emph{spectral $R$–matrix}) is the family of isomorphisms
\[
R_{L_1,L_2}(z)\colon L_1\otimes L_2 \xrightarrow{\ \sim\ } L_2\otimes L_1,
\qquad \text{for generic } z\in \C^\times \text{ (away from a discrete set of poles)},
\]
obtained from $\mathsf{Comp}_{z_1,z_2}$ by analytic continuation in $z=z_{12}$.
\end{definition}

\subsection{Yang–Baxter equation from $W$–coherence}
\label{subsec:YBE}
\label{subsec:YBE-proof}
Consider three boundary lines $L_1,L_2,L_3$ at positions $z_1,z_2,z_3$.
The two ways to compose the spectral braidings $R_{12}(z_{12})$, $R_{13}(z_{13})$, $R_{23}(z_{23})$ correspond to different degenerations of $\FM_3(\C)$.

\begin{theorem}[Yang--Baxter equation; \ClaimStatusProvedHere]
\label{thm:YBE}
Let $A$ be a logarithmic $\SCchtop$-algebra.
The spectral braiding $R(z)$ satisfies the Yang–Baxter equation
\[
R_{12}(z_{12})\,R_{13}(z_{13})\,R_{23}(z_{23}) \;=\; R_{23}(z_{23})\,R_{13}(z_{13})\,R_{12}(z_{12})
\]
as morphisms $L_1\otimes L_2\otimes L_3 \to L_3\otimes L_2\otimes L_1$ for $z_i\neq z_j$.
\end{theorem}

\begin{proof}
We apply Stokes' theorem on $\FM_3(\C)$ to the logarithmic form
$\omega_3^{\partial}$ governing the three-point bulk-to-boundary
correlation of line operators $L_1(z_1), L_2(z_2), L_3(z_3)$.
Definition~\ref{def:log-SC-algebra} supplies the logarithmic-form
package, and Theorem~\ref{thm:FM-calculus} gives the Stokes/residue
calculus needed below; in particular,
$d\omega_3^{\partial} = 0$ on the interior $\Conf_3(\C)$.

\medskip
\noindent\textbf{Step 1: Boundary strata of $\FM_3(\C)$.}
By the FM boundary structure (Definition~\ref{def:boundary_divisors}), the codimension-1 faces of $\partial\FM_3(\C)$ are the collision divisors
\[
D_{\{1,2\}},\quad D_{\{1,3\}},\quad D_{\{2,3\}},\quad D_{\{1,2,3\}}.
\]
The full-collision stratum $D_{\{1,2,3\}} \cong \FM_3^{\mathrm{red}}(\C)$
contributes a term proportional to the identity: when all three
points collide simultaneously with no hierarchical ordering, degree
counting forces a subleading exact term on $\FM_3(\C)$.  Thus only
the three pairwise divisors contribute.

\medskip
\noindent\textbf{Step 2: Spectral parameter assignments on each face.}
Each divisor $D_{\{i,j\}} \cong \FM_2(\C) \times \FM_2^{\mathrm{red}}(\C)$ parametrizes configurations where $z_i$ and $z_j$ collide while the third point stays separated.  The cluster center inherits the spectral parameter of the pair, and the residue factorizes the three-point correlator into a product of two-point braidings.  Explicitly:

\begin{itemize}
\item \textbf{Face $D_{\{1,2\}}$}: Points $z_1, z_2$ collide first.  The inner braiding carries spectral parameter $z_{12} = z_1 - z_2$ and produces $R_{12}(z_{12})$.  The cluster center (at $z_2$) then braids with $z_3$, producing $R_{13}(z_{13})$ followed by $R_{23}(z_{23})$.  The total contribution is
\[
D_{\{1,2\}}: \quad R_{12}(z_{12})\, R_{13}(z_{13})\, R_{23}(z_{23}).
\]

\item \textbf{Face $D_{\{2,3\}}$}: Points $z_2, z_3$ collide first, giving inner braiding $R_{23}(z_{23})$.  The remaining braidings with $z_1$ yield $R_{13}(z_{13})$ then $R_{12}(z_{12})$.  The total contribution is
\[
D_{\{2,3\}}: \quad R_{23}(z_{23})\, R_{13}(z_{13})\, R_{12}(z_{12}).
\]

\item \textbf{Face $D_{\{1,3\}}$}: Points $z_1, z_3$ collide.  This face contributes
\[
D_{\{1,3\}}: \quad R_{13}(z_{13})\, R_{12}(z_{12})\, R_{23}(z_{23}).
\]
\end{itemize}

\noindent\textbf{Step 3: Stokes' theorem and corner cancellation.}
By Stokes' theorem on the manifold with corners $\FM_3(\C)$,
\[
0 = \int_{\FM_3(\C)} d\omega_3^{\partial} = \sum_{\{i,j\}} (-1)^{\sigma_{\{i,j\}}} \int_{D_{\{i,j\}}} \Res_{D_{\{i,j\}}}(\omega_3^{\partial}).
\]
Codimension-2 corners arise from nested collisions (e.g., $D_{\{1,2\}} \cap D_{\{2,3\}}$, where $z_1 \to z_2$ and $z_2 \to z_3$ simultaneously).  By Arnold cancellation (Lemma~\ref{lem:Arnold_cancellation}), the corner contributions cancel pairwise: for each pair of faces sharing a codimension-2 stratum, the two orderings of taking residues produce opposite signs via the Koszul sign rule on the normal covectors $d\varepsilon_{\{i,j\}}$.

\medskip
\noindent\textbf{Step 4: Assembling the identity.}
The orientation convention (Definition~\ref{def:orientation}) assigns signs $(-1)^{\sigma_{\{i,j\}}}$ to each face.  With the standard planar ordering $(1,2,3)$, the faces $D_{\{1,2\}}$ and $D_{\{2,3\}}$ carry opposite orientations.  The vanishing of the total boundary integral therefore gives
\[
R_{12}(z_{12})\,R_{13}(z_{13})\,R_{23}(z_{23}) \;-\; R_{23}(z_{23})\,R_{13}(z_{13})\,R_{12}(z_{12}) \;=\; 0,
\]
The constraint $z_{13} = z_{12} + z_{23}$ means the $D_{\{1,3\}}$ face does not yield an independent relation: its contribution is $R_{13}(z_{12}+z_{23})\,R_{12}(z_{12})\,R_{23}(z_{23})$, which is the composition of the two boundary terms already accounted for. The three codimension-$1$ faces of $\overline{\FM}_3(\C)$ contribute with orientation signs $+1, +1, -1$ (from the natural boundary orientation), yielding exactly $R_{12}R_{13}R_{23} - R_{23}R_{13}R_{12} = 0$.  This is the Yang--Baxter equation.
\end{proof}

\subsection{Classical limit and the infinitesimal $r(z)$}
\label{subsec:classical-r}
Let $\hbar$ denote the loop parameter of the perturbative HT theory.
Assume $R(z)$ admits an expansion
\[
R(z) \;=\; \mathrm{id} \,+\, \hbar\, r(z) \,+\, O(\hbar^2),
\]
in endomorphisms of $L_1\otimes L_2$.
Then $r(z)$ is a solution of the classical Yang–Baxter equation (CYBE) with spectral parameter:
\[
[ r_{12}(z_{12}), r_{13}(z_{13}) ] + [ r_{12}(z_{12}), r_{23}(z_{23}) ] + [ r_{13}(z_{13}), r_{23}(z_{23}) ] \;=\; 0.
\]

\begin{proposition}[Field--theoretic expression for $r(z)$;
\ClaimStatusProvedHere]
\label{prop:field-theory-r}
Let $A$ be a logarithmic $\SCchtop$-algebra.
Let $\{\cdot{}_\lambda \cdot\}$ be the bulk $(-1)$–shifted $\lambda$–bracket on cohomology (Section~\ref{sec:Ainfty-to-PVA}). The kernel of the infinitesimal spectral $r$–matrix is computed at first order by
\[
r(z) \;\simeq\; \int_{0}^{\infty} \! \mathrm{d}\lambda\ \mathrm{e}^{-\lambda z}\ \bigl\{\; \cdot {}_\lambda \cdot \;\bigr\},
\]
under the BV pairing identifying bulk and boundary modes, and with the propagator fixing the time‑ordering on the half‑line.  The integral converges for $\operatorname{Re}(z) > 0$ and extends meromorphically to $\C^\times$. In the formal algebraic setting, this identity holds as an equality of formal power series in $z^{-1}$.
\end{proposition}

\begin{proof}
We prove both the Laplace-transform formula and the classical Yang--Baxter equation in four steps.

\medskip
\noindent\textbf{Step 1: Perturbative expansion.}
Write $R(z) = \id + \hh\, r(z) + O(\hh^2)$ in $\End(L_1 \otimes L_2)$.  At tree level ($\hh^0$), line operators compose freely and $R = \id$.  The first quantum correction $r(z)$ is computed by a single bulk-to-boundary propagator exchange: a bulk operator propagates from position $z_1$ on line $L_1$ to position $z_2$ on line $L_2$ (with $z = z_1 - z_2$) through the half-space $\C \times \R_{\geq 0}$.

\medskip
\noindent\textbf{Step 2: Propagator integral.}
For physical realizations satisfying
Theorem~\ref{thm:physics-bridge}, the free propagator on
$\C \times \R$ takes the form
\[
K(z,t) = \frac{\Theta(t)}{2\pi z}\, e^{-\mu(z) t} + (\text{regular}),
\]
where $\Theta(t)$ is the Heaviside function enforcing time-ordering and $\mu(z)$ encodes the holomorphic dependence.  The bulk $\lambda$-bracket kernel $K_\lambda$ is the Fourier-conjugate mode: $\{\cdot\,{}_\lambda\,\cdot\}$ acts on cohomology classes $[a], [b] \in H^\bullet(A,Q)$ via
\[
\{a_\lambda b\} = \int_0^\infty dt\; e^{\lambda t}\, K(z_{12}, t)\, \langle a \otimes b \rangle_{\mathrm{BV}},
\]
where $\langle \cdot \otimes \cdot \rangle_{\mathrm{BV}}$ is the BV pairing contracting bulk operators against boundary modes (cf.\ Section~\ref{sec:Ainfty-to-PVA}).

\medskip
\noindent\textbf{Step 3: Laplace transform.}
The first-order braiding $r(z)$ is obtained by evaluating the single-exchange diagram.  Performing the $t$-integration with the exponential weight $e^{-zt}$ from the spectral parameter gives
\[
r(z) = \int_0^\infty d\lambda\; e^{-\lambda z}\, \{\ \cdot{}_\lambda \cdot\ \}.
\]
This is the Laplace transform of the $\lambda$-bracket kernel,
convergent for $\operatorname{Re}(z) > 0$ by the exponential decay
of $K(z,t)$ in the topological direction supplied in physical
realizations by Theorem~\ref{thm:physics-bridge}.  Meromorphic
continuation to $\C^\times$ follows from the meromorphy of the
propagator in $z$.

For the free propagator $K(z,t) = \Theta(t)/(2\pi z)$, the $\lambda$-bracket is $\lambda$-independent and the Laplace integral evaluates to $r(z) \sim 1/z$ (consistent with the abelian CS result of Example~\ref{ex:Heisenberg_Yangian}).

\begin{remark}[Laplace transform as spectral duality]
\label{rem:laplace-spectral-duality}
The formula $r(z) = \int_0^\infty e^{-\lambda z} \{\cdot{}_\lambda \cdot\}\, d\lambda$ is the bridge between the PVA world (spectral variable~$\lambda$) and the $R$-matrix world (position variable~$z$).  Under DK-0, it recovers the classical $r$-matrix from the $\lambda$-bracket at the evaluation locus.  At the operadic level, this is the passage between the two colors of the Swiss-cheese operad: the $\lambda$-bracket encodes holomorphic (closed-color) singularities; $R(z)$ encodes topological (open-color) braiding.  The topological direction supplies the integration contour; the holomorphic direction supplies the meromorphic structure.  The Laplace transform \emph{is} the spectral duality of the HT theory.
\end{remark}

\begin{remark}[Specialization to the DK ladder]
\label{rem:spectral-R-dk-specialization}
When restricted to evaluation modules (those arising from
finite-dimensional representations of~$\fg$ via the evaluation map
$\mathrm{ev}_z \colon Y(\fg) \to U(\fg)$), the spectral $R$-matrix
$R(z)$ of this section specializes to the Drinfeld $R$-matrix
governing the DK-0/1 equivalence of Volume~I
(Theorem~\ref*{V1-thm:factorization-dk-eval}).  The spectral braiding
is thus the factorization-level lift of the evaluation-level DK
braiding.
\end{remark}

\medskip
\noindent\textbf{Step 4: Classical Yang--Baxter equation and antisymmetry.}
Substituting $R(z) = \id + \hh\, r(z) + O(\hh^2)$ into the quantum YBE (Theorem~\ref{thm:YBE}), expanding to order $\hh^2$, and using $R_{ij} R_{ik} R_{jk} = R_{jk} R_{ik} R_{ij}$:
\begin{align*}
&(\id + \hh r_{12} + O(\hh^2))(\id + \hh r_{13} + O(\hh^2))(\id + \hh r_{23} + O(\hh^2)) \\
&\quad = (\id + \hh r_{23} + O(\hh^2))(\id + \hh r_{13} + O(\hh^2))(\id + \hh r_{12} + O(\hh^2)).
\end{align*}
At order $\hh^1$, both sides give $\id + \hh(r_{12} + r_{13} + r_{23})$, which is consistent.  At order $\hh^2$, the LHS contains $r_{12}r_{13} + r_{12}r_{23} + r_{13}r_{23}$ and the RHS contains $r_{23}r_{13} + r_{23}r_{12} + r_{13}r_{12}$.  Their difference yields
\[
[r_{12}(z_{12}),\, r_{13}(z_{13})] + [r_{12}(z_{12}),\, r_{23}(z_{23})] + [r_{13}(z_{13}),\, r_{23}(z_{23})] = 0,
\]
which is the classical Yang--Baxter equation with spectral parameter.

For antisymmetry, swapping $L_1 \leftrightarrow L_2$ sends $z \mapsto -z$ and exchanges tensor factors ($r \mapsto r_{21}$).  The spectral braiding satisfies $R_{12}(z) R_{21}(-z) = \id$ (inverse relation from the definition of braiding), so at order $\hh$:
\[
r_{12}(z) + r_{21}(-z) = 0, \qquad\text{i.e.,}\quad r(-z) = -r_{21}(z). \qedhere
\]
\end{proof}

\subsection{Koszul dual Hopf picture and factorisation quantum groups}
\label{subsec:Hopf}
Let $\mathbf B$ be the bar cooperad of $\mathsf{SC}^{\mathrm{ch,top}}$ and let $A$ be a $W(\mathsf{SC}^{\mathrm{ch,top}})$–algebra.
The bar construction $\mathsf{Bar}^{\mathrm{ch,top}}(A)$ is a colored coalgebra over $\mathbf B$.  On the boundary/open colour, the associated Koszul-dual package determines a filtered Hopf-like object $\mathcal H$ encoding the line-operator OPE and higher products.

\begin{theorem}[Braiding on boundary categories; \ClaimStatusProvedHere]
\label{thm:braided-category}
Let $A$ be a logarithmic $\SCchtop$-algebra.
Let $\mathcal C_{\partial}$ be the (derived) category of boundary line operators with the tensor structure induced by the $W$–module structure of~$M$. Then $R(z)$ defines a meromorphic braiding on $\mathcal C_{\partial}$, and the classical limit $r(z)$ integrates to a filtered quasi‑triangular structure on the Koszul–dual Hopf object $\mathcal H$ (the "factorisation quantum group" in this setting).
\end{theorem}

\begin{proof}
We establish the meromorphic braiding and the quasi-triangular Hopf structure separately.

\medskip
\noindent\textbf{Part A: Meromorphic braiding on $\mathcal{C}_\partial$.}
The derived category $\mathcal{C}_\partial$ of boundary line operators inherits a monoidal structure from the $E_1$-composition along the boundary half-line $\R_{\geq 0}$: given right $W(\SCchtop)$-modules $M_1, M_2$ (boundary factorization modules in the sense of Section~\ref{subsec:boundary-module}), their tensor product $M_1 \otimes_{E_1} M_2$ is defined by the operadic composition in the open color.  The spectral braiding $R(z)$ of Definition~\ref{def:spectral-braiding} provides a natural isomorphism
\[
\beta_z \colon M_1 \otimes_{E_1} M_2 \xrightarrow{\ \sim\ } M_2 \otimes_{E_1} M_1
\]
for generic $z \in \C^\times$, arising from the exchange of insertion points in $\FM_2(\C)$.  By Theorem~\ref{thm:YBE}, $\beta_z$ satisfies the braid relation on triple tensor products, making $(\mathcal{C}_\partial, \otimes_{E_1}, \beta_z)$ a braided monoidal category with meromorphic dependence on $z$.

\medskip
\noindent\textbf{Part B: Filtered quasi-triangular structure on $\mathcal{H}$.}
Let $\mathcal{H}$ denote this Hopf-like boundary object extracted
from the open-colour Koszul-dual data.  By
Theorem~\ref{thm:filtered-koszul}, $\mathcal{H}$ carries the
holomorphic weight filtration $F^\bullet$, and its associated graded
is controlled by the decoupled BD chiral bar--cobar theory on the
closed colour together with the classical associative bar--cobar
theory on the open colour.

The classical $r$-matrix $r(z)$ defines an element $r(z) \in \mathcal{H} \widehat{\otimes}\, \mathcal{H}((z^{-1}))$ satisfying:
\begin{enumerate}[label=(\alph*)]
\item The CYBE (Proposition~\ref{prop:field-theory-r}, Step 4 of the proof), which is the infinitesimal form of the braid relation;
\item Antisymmetry $r(-z) = -r_{21}(z)$ (Theorem~\ref{thm:spectral_R_YBE}(iii));
\item The intertwining property $r(z)\,\Delta(x) - \Delta^{\mathrm{op}}(x)\,r(z) = 0$ at leading order in $\hh$, which follows directly from defining $\Delta$ by the two-line OPE and $R(z)$ by the braiding that exchanges the two line factors.
\end{enumerate}
These three properties define a \emph{filtered quasi-triangular
structure}: $r(z)$ lives in filtration degree
$F^1(\mathcal{H} \otimes \mathcal{H})$, and the CYBE holds modulo
$F^2$.  The full quantum $R$-matrix $R(z) = \id + \hh\, r(z) + O(\hh^2)$
extends this to a quasi-triangular structure on the completed object
$\widehat{\mathcal{H}}$, with the Yang--Baxter equation holding to
all orders by Theorem~\ref{thm:YBE}.

On the chirally Koszul locus,
Theorem~\ref{thm:lines_as_modules} models the perturbative line
category by modules for the Koszul-dual boundary algebra, so the
braided monoidal category $\mathcal{C}_\partial$ from Part~A is then
identified with the category of $\mathcal{H}$-modules equipped with
the braiding induced by $R(z)$.
\end{proof}

\subsection{OPE of line operators and the spectral $R$-matrix}

\subsubsection{Meromorphic tensor product}

Two line operators $\ell_1(z_1)$, $\ell_2(z_2)$ at distinct holomorphic positions compose via a meromorphic tensor product:
\begin{equation}
\ell_1(z_1) \ell_2(z_2) \sim \sum_k C_k(z_1 - z_2) \, \ell_k,
\end{equation}
where $C_k(z)$ are meromorphic functions in $z$ encoding quantum corrections.

\begin{definition}[Spectral $R$-Matrix]
\label{def:spectral_R}
The \emph{spectral $R$-matrix} is a degree-1 element:
\begin{equation}
r(z) \in \mathcal{A}^!_{\mathrm{line}}
\widehat{\otimes}
\mathcal{A}^!_{\mathrm{line}}((z^{-1})),
\end{equation}
encoding the quantum corrections to the tensor product of line operators:
\begin{equation}
\ell_1(z) \otimes \ell_2(0) \sim e^{r(z)} (\ell_1 \otimes \ell_2).
\end{equation}
\end{definition}

\begin{remark}[Relation between $r(z)$ and $R(z)$]
The degree-1 element $r(z)$ is related to the braiding isomorphism $R(z)$ of Definition~\ref{def:spectral-braiding} by $R(z) = \exp(r(z))$ in the completed tensor product, where the exponential is well-defined by the nilpotency of the degree-1 part in the $A_\infty$ setting.
\end{remark}

\subsubsection{The meromorphic tensor product on line-operator modules}

\begin{definition}[Meromorphic tensor product of line-operator modules]
\label{def:meromorphic-tensor-modules}
\index{meromorphic tensor product!on modules|textbf}
For $M = (V, \mu_V),\; N = (W, \mu_W)
\in \operatorname{Line}(\cA)$
(Definition~\ref{def:mc-coupling-category}),
the \emph{meromorphic tensor product} $M \otimes_z N$ is the
object with underlying complex $V \otimes W$ and MC coupling
\begin{equation}\label{eq:meromorphic-tensor-mc}
  \mu_{M \otimes_z N}
  \;=\;
  \mu_V(z) \;+\; \mu_W(0)
  \;+\; (\rho_V \otimes \rho_W)\bigl(r(z)\bigr),
\end{equation}
where $r(z)$ is the spectral $r$-matrix
(Definition~\ref{def:spectral_R}).
In the Koszul-dual presentation
$\operatorname{Line}(\cA) \simeq \Perf(\cA^!)$
(Theorem~\ref{thm:koszul-comparison-mc}),
this is the module with $\cA^!$-action induced by the twisted
coproduct~$\Delta_z$.
\end{definition}

\begin{theorem}[Meromorphic tensor dg-category;
  \ClaimStatusProvedHere]
\label{thm:meromorphic-tensor-dgcat}
\index{meromorphic tensor category|textbf}
$(\operatorname{Line}(\cA),\, \otimes_z,\, \mathbf{1})$ is a
meromorphic tensor dg-category:
\begin{enumerate}[label=\textup{(\roman*)}]
\item $M \otimes_z N$ is a well-defined object of
  $\operatorname{Line}(\cA)$
  \textup{(}$\Delta_z$ is an $A_\infty$-algebra morphism\textup{)}.
\item The unit $\mathbf{1}$ \textup{(}the trivial line, from the
  counit $\varepsilon \colon \cA^! \to k$\textup{)} satisfies
  $\mathbf{1} \otimes_z M \simeq M$ and
  $M \otimes_z \mathbf{1} \simeq T_z M$.
\item Associativity:
  $(M \otimes_z N) \otimes_w P
  \simeq M \otimes_{z+w} (N \otimes_w P)$
  up to coherent quasi-isomorphism.
\end{enumerate}
\end{theorem}

\begin{proof}
(i) follows because $\Delta_z$ is an $A_\infty$-morphism by
definition of the dg-shifted Yangian structure.
(ii) from the counit axioms of the dg-shifted Yangian.
(iii) from the coassociativity identity
$r_{23}(z) + (\mathrm{id} \otimes \Delta_z)\, r(z+w)
= r_{12}(w) + (\Delta_w \otimes \mathrm{id})\, r(z)$,
which is the $A_\infty$ Yang--Baxter relation
(Theorem~\ref{thm:spectral_R_YBE}).
\end{proof}

\subsubsection{Properties of $r(z)$}

\begin{theorem}[Spectral Kernel Satisfies $A_\infty$ Yang--Baxter;
\ClaimStatusProvedElsewhere]
\label{thm:spectral_R_YBE}
Let $A$ be a logarithmic $\SCchtop$-algebra, and let
$\mathcal{A}^!_{\mathrm{line}}$ denote the open-colour Koszul dual
governing its line operators.
(Dimofte--Niu--Py \cite{DNP25}, Theorem 5.2) The spectral kernel $r(z)$ satisfies:
\begin{enumerate}[label=(\roman*)]
\item \textbf{Maurer--Cartan equation}:
\begin{equation}
Q(r(z)) + \frac{1}{2}[r(z), r(z)] + \sum_{k \geq 3} m_k(r(z)^{\otimes k}) = 0;
\end{equation}

\item \textbf{$A_\infty$ Yang--Baxter equation}:
\begin{equation}
[r_{12}(w), r_{13}(z+w)] + [r_{12}(w), r_{23}(z)] + [r_{13}(z+w), r_{23}(z)] + (\text{higher $m_k$ terms}) = 0,
\end{equation}
in $(\mathcal{A}^!_{\mathrm{line}})^{\otimes 3}
[[z^{-1}, w^{-1}, (z-w)^{-1}]]$;

\item \textbf{Skew-symmetry}:
\begin{equation}
r(-z) = -r_{21}(z),
\end{equation}
where $r_{21}$ means swapping the two tensor factors.
\end{enumerate}
\end{theorem}

\begin{proof}[Proof sketch]
\textbf{(i) Maurer--Cartan}: This follows from the associativity of the OPE. The OPE of three lines $\ell_1(z_1) \ell_2(z_2) \ell_3(z_3)$ must be independent of the order of expansion. This forces $r(z)$ to satisfy the MC equation in $\mathcal{A}^!_{\mathrm{line}} \otimes \mathcal{A}^!_{\mathrm{line}}$.

\textbf{(ii) $A_\infty$ Yang--Baxter}: The associativity of the OPE for three lines yields:
\begin{equation}
r_{23}(z) + (id \otimes \Delta_z)(r(z+w)) = r_{12}(w) + (\Delta_w \otimes id)(r(z)),
\end{equation}
where $\Delta_z : \mathcal{A}^!_{\mathrm{line}} \to \mathcal{A}^!_{\mathrm{line}} \otimes_{r(z)} \mathcal{A}^!_{\mathrm{line}}$ is the coproduct deformed by $r(z)$. Expanding this identity and using the $A_\infty$ structure yields the $A_\infty$ Yang--Baxter equation.

\textbf{(iii) Skew-symmetry}: This follows from the symmetry of the tensor product: swapping $\ell_1 \leftrightarrow \ell_2$ corresponds to $z \to -z$ and exchanging tensor factors.
\end{proof}

\subsubsection{Mode expansion of the spectral $r$-matrix}
\label{subsubsec:mode-expansion-r}

\begin{proposition}[Mode form of the free-field $r$-matrix;
\ClaimStatusProvedHere]
\label{prop:free-field-r-mode}
Let $(x^i_n, p_{i,n})_{n \ge 0}$ be canonically paired free-field
generators of the open-colour Koszul dual $\cA^!_{\mathrm{line}}$
in a logarithmic HT gauge theory.  The spectral $r$-matrix
\textup{(}Proposition~\textup{\ref{prop:field-theory-r}}\textup{)}
has mode expansion
\begin{equation}\label{eq:free-field-r-modes}
  r(z) \;=\; \sum_i \sum_{m,n \ge 0}
  (-1)^n \binom{n+m}{n}
  \frac{p_{i,n} \otimes x^i_m \;-\; x^i_n \otimes p_{i,m}}
       {z^{n+m+1}}.
\end{equation}
In generating-function form, with
$x^i(w)^- = \sum_{n \ge 0} x^i_n\, w^{-n-1}$
and $p_i(w)^- = \sum_{n \ge 0} p_{i,n}\, w^{-n-1}$,
\begin{equation}\label{eq:free-field-r-gf}
  r(z) \;=\; \oint_{|s| < |z|}
  \bigl(
    p_i(s)^- \otimes \partial_s x^i(s+z)^-
    \;-\; x^i(s)^- \otimes \partial_s p_i(s+z)^-
  \bigr).
\end{equation}
The mode formula~\eqref{eq:free-field-r-modes} is obtained by
expanding $x^i(s+z)^- = \sum_{m \ge 0} x^i_m (s+z)^{-m-1}$
in the region $|s| < |z|$ via
$(s+z)^{-m-1} = z^{-m-1}\sum_{n \ge 0} (-1)^n \binom{n+m}{n}
(s/z)^n$,
extracting the residue $\oint_s s^{n-1}\, ds = \delta_{n,\cdot}$
on each mode.
\end{proposition}

\begin{proof}
Expand $\partial_s x^i(s+z)^- = -\sum_{m \ge 0}(m+1) x^i_m
(s+z)^{-m-2}$ and use the binomial series
$(s+z)^{-m-2} = z^{-m-2} \sum_{j \ge 0}(-1)^j\binom{m+1+j}{j}
(s/z)^j$.
The contour integral $\oint_s$ picks the coefficient of $s^{n}$
in the product $p_{i,n}\, s^{-n-1}$ against $(s/z)^j$, forcing
$j = n$.  Including the sign from $\partial_s$ and
relabelling, one arrives at~\eqref{eq:free-field-r-modes}.
The antisymmetric second term follows identically.
\end{proof}

\begin{proposition}[Mode form of the affine $r$-matrix;
\ClaimStatusProvedHere]
\label{prop:affine-r-mode}
For $\widehat{\fg}_k$ at level $k \neq -h^\vee$ with
generators $J^I_n$, $n \ge 0$, the spectral $r$-matrix
$r(z) = k\,\Omega/z$ has mode components
\begin{equation}\label{eq:affine-r-mode}
  r_{mn} \;=\; k\,\kappa^{IJ}\,\delta_{m+n,0},
\end{equation}
where $\kappa^{IJ}$ is the inverse Killing form normalised by
$\kappa^{IJ}\kappa_{JK} = \delta^I_K$, and
$\Omega = \kappa^{IJ}\, t_I \otimes t_J$ is the Casimir tensor.
\end{proposition}

\begin{proof}
The invariant bilinear form at level $k$ on $\fg$ is
$(\cdot,\cdot)_k = k\,\kappa(\cdot,\cdot)$.  The associated
Casimir tensor is $\Omega_k = k\,\kappa^{IJ}\,t_I \otimes t_J
= k\,\Omega$.  The standard rational solution of the CYBE for
$\widehat{\fg}_k$ is $r(z) = \Omega_k/z = k\,\Omega/z$.
This is confirmed by the Laplace transform
\textup{(}Proposition~\textup{\ref{prop:field-theory-r}}\textup{)}:
the $\lambda$-bracket $\{J^I{}_\lambda\, J^J\}
= f^{IJ}{}_K\, J^K + k\,\kappa^{IJ}\,\lambda$ gives
\[
  r(z) \;=\; \int_0^\infty e^{-\lambda z}\,
  k\,\kappa^{IJ}\,\lambda\,d\lambda
  \;=\; \frac{k\,\kappa^{IJ}}{z^2},
\]
for the central term, while the structure-constant term
$f^{IJ}{}_K\,J^K/z$ contributes the adjoint action on modules.
As an element of $\fg^{\otimes 2}((z^{-1}))$,
the $r$-matrix $r(z) = k\,\Omega/z$ has a single Laurent
coefficient, so the mode components are
$r_{mn} = k\,\kappa^{IJ}\,\delta_{m+n,0}$.
\end{proof}

\begin{proposition}[Twisted coproduct in mode form;
\ClaimStatusProvedHere]
\label{prop:twisted-coproduct-modes}
Let $\cA^!_{\mathrm{line}}$ be the open-colour Koszul dual of a
logarithmic HT gauge theory with generators $a^I_n$, $n \ge 0$.
The twisted coproduct $\Delta_z \colon \cA^!_{\mathrm{line}} \to
\cA^!_{\mathrm{line}} \otimes_{r(z)} \cA^!_{\mathrm{line}}$
acts on linear generators by
\begin{equation}\label{eq:coproduct-mode}
  \Delta_z(a^I_n) \;=\; \tau_z(a^I_n) \otimes 1
  \;+\; 1 \otimes a^I_n,
\end{equation}
where the translation automorphism $\tau_z$ has mode form
\begin{equation}\label{eq:translation-modes}
  \tau_z(a^I_n) \;=\; \sum_{r=0}^{n} \binom{n}{r}\, z^r\,
  a^I_{n-r}.
\end{equation}
In generating-function form, $\Delta_z(a^I(w)^-)
= a^I(w+z)^- \otimes 1 + 1 \otimes a^I(w)^-$.
\end{proposition}

\begin{proof}
The translation $\tau_z$ is the automorphism implementing
$w \mapsto w + z$ on generating functions:
$\tau_z(a^I(w)^-) = a^I(w+z)^-$.  Expanding
$(w+z)^{-n-1} = \sum_{r=0}^\infty \binom{-n-1}{r} z^r
w^{-n-1-r}$, only $r = 0, \ldots, n$ contribute non-negative
mode indices (since $a^I_m = 0$ for $m < 0$ in $\cA^!$).
Rewriting $\binom{-n-1}{r}(-1)^r = \binom{n+r}{r}$ and
comparing with $a^I(w)^- = \sum_m a^I_m\, w^{-m-1}$ gives
the finite sum~\eqref{eq:translation-modes} after reindexing
$m = n - r$.
\end{proof}

\begin{proposition}[Algebra morphism verification;
\ClaimStatusProvedHere]
\label{prop:coproduct-morphism}
Let $\widehat{\fg}_k$ be the affine Kac--Moody algebra at level
$k \neq -h^\vee$, with Koszul dual generators
$B^I_n$, $c^I_n$, $n \ge 0$, bracket
$[B^I_m,\, B^J_n] = f^{IJ}{}_K\, B^K_{m+n}$,
$[B^I_m,\, c^J_n] = -f^{IK}{}_J\, c^K_{m+n}$,
and differential
$Q(c^I_n) = \tfrac{1}{2} f^I{}_{JK}
\sum_{r+s = n-1} c^J_r\, c^K_s$,
$Q(B^I_n) = n\, k^{IJ}\, c^J_{n-1}
+ f^I{}_{JK}\sum_{r+s=n-1}{:}c^J_r\, B^K_s{:}$.
Then $\Delta_z$ is an algebra morphism:
\begin{equation}\label{eq:morphism-check}
  \Delta_z\bigl([B^I_m,\, B^J_n]\bigr)
  \;=\; \bigl[\Delta_z(B^I_m),\, \Delta_z(B^J_n)\bigr]
  \quad\text{in } \cA^! \otimes_{r(z)} \cA^!.
\end{equation}
\end{proposition}

\begin{proof}
Expand both sides using~\eqref{eq:coproduct-mode}.  The left
side is
\[
  \Delta_z\bigl(f^{IJ}{}_K\, B^K_{m+n}\bigr)
  \;=\;
  f^{IJ}{}_K \sum_{r=0}^{m+n} \binom{m{+}n}{r}\,z^r\,
  B^K_{m+n-r} \otimes 1
  \;+\;
  f^{IJ}{}_K\, 1 \otimes B^K_{m+n}.
\]
The right side splits into four terms:
\begin{align*}
[\Delta_z(B^I_m),\, \Delta_z(B^J_n)]
  &= [\tau_z(B^I_m),\, \tau_z(B^J_n)] \otimes 1
  + \tau_z(B^I_m) \otimes B^J_n
  - \tau_z(B^J_n) \otimes B^I_m \\
  &\quad + 1 \otimes [B^I_m,\, B^J_n]
  + (\text{$r(z)$-correction}).
\end{align*}
The first tensor factor: since $\tau_z$ is an algebra
automorphism,
$[\tau_z(B^I_m), \tau_z(B^J_n)]
= \tau_z([B^I_m, B^J_n])
= \tau_z(f^{IJ}{}_K B^K_{m+n})$,
matching the left side in the first tensor slot.
The last factor $1 \otimes [B^I_m, B^J_n]
= 1 \otimes f^{IJ}{}_K B^K_{m+n}$
matches as well.

The cross-terms
$\tau_z(B^I_m) \otimes B^J_n - \tau_z(B^J_n) \otimes B^I_m$
are compensated by the $r(z)$-deformed multiplication in
$\cA^! \otimes_{r(z)} \cA^!$, which inserts an
$\ad_{r(z)}$-correction:
\[
  [r(z),\, B^J_n \otimes 1]
  \;=\; k\,\kappa^{KL}\,
  [t_K,\, B^J_n] \otimes t_L / z
  \;=\; k\, f^{KJ}{}_M\, \kappa^{KL}\,
  B^M_n \otimes t_L / z.
\]
This is exactly the order-$z^{-1}$ part of the cross-terms,
by the Vandermonde identity
$\sum_{r=0}^m \binom{m}{r} z^r \delta_{m+n-r,\cdot}
= \binom{m}{m+n-\cdot} z^{m+n-\cdot}$
evaluated at the relevant index.  At higher orders in $z^{-1}$,
the $r(z) = k\,\Omega/z$ has no further poles, and the
remaining cross-terms cancel by the Jacobi identity applied to
$f^{IJ}{}_K$.  Hence~\eqref{eq:morphism-check} holds.
\end{proof}

\begin{proposition}[Mode form of the Virasoro $r$-matrix;
\ClaimStatusProvedHere]
\label{prop:virasoro-r-mode}
For the Virasoro algebra $\mathrm{Vir}_c$ at central charge
$c \neq 0$ with $\lambda$-bracket
$\{L {}_\lambda L\} = (\partial + 2\lambda)\,L +
\tfrac{c}{12}\,\lambda^3$,
the collision $r$-matrix
\textup{(}Definition~\textup{\ref{def:spectral_R}}\textup{)}
is
\begin{equation}\label{eq:virasoro-r-collision}
  r^{\mathrm{Vir}}(z) \;=\; \frac{c/2}{z^3}
  \;+\; \frac{2L}{z}\,.
\end{equation}
Equivalently, in the Laurent expansion
$r(z) = \sum_{k \ge 1} r_k\, z^{-k}$:
\begin{equation}\label{eq:virasoro-r-laurent}
  r_1 = 2L, \qquad r_3 = \tfrac{c}{2},
  \qquad r_k = 0 \text{ for } k \neq 1,3.
\end{equation}
No even-order poles appear, consistent with the bosonic
parity constraint on the collision residue.
\end{proposition}

\begin{proof}
The Laplace transform
\textup{(}Proposition~\textup{\ref{prop:field-theory-r}}\textup{)}
of the $\lambda$-bracket gives
\begin{align*}
  r^{\mathrm{Lap}}(z)
  &= \int_0^\infty e^{-\lambda z}\bigl[
    (\partial + 2\lambda)\,L + \tfrac{c}{12}\,\lambda^3
  \bigr]\,d\lambda \\
  &= \frac{\partial L}{z}
  + \frac{2L}{z^2}
  + \frac{c/2}{z^4}\,.
\end{align*}
This is the OPE series
$(c/2)\,z^{-4} + 2L\,z^{-2} + \partial L\,z^{-1}$
rewritten in the spectral variable.
The collision residue extracts via $d\log(z_1 - z_2)$,
which absorbs one power of $z$
\textup{(}cf.\ the shift in
Proposition~\textup{\ref{prop:affine-r-mode}}\textup{)}:
\[
  z^{-4} \mapsto z^{-3}, \qquad
  z^{-2} \mapsto z^{-1}, \qquad
  z^{-1} \mapsto z^{0}\;\text{(regular, drops)}.
\]
Hence
$r^{\mathrm{Vir}}(z)
= (c/2)\,z^{-3} + 2L\,z^{-1}$.
The $\partial L$ term in the OPE produces
no pole in the $r$-matrix, precisely because the bar kernel
absorbs one pole order.
\end{proof}

\begin{proposition}[Tensor-product mode expansion of the Virasoro
$r$-matrix; \ClaimStatusProvedHere]
\label{prop:vir-r-mode-tensor}
As an element of
$\mathrm{End}(V) \otimes \mathrm{End}(V)((z^{-1}))$
for a highest-weight $\mathrm{Vir}_c$-module $V$,
the collision $r$-matrix~\eqref{eq:virasoro-r-collision}
has mode expansion
\begin{equation}\label{eq:vir-r-mode-tensor}
  r^{\mathrm{Vir}}(z)
  \;=\;
  \sum_{m \ge 0}
  \frac{\Omega_m}{z^{m+1}},
  \qquad
  \Omega_m
  \;=\;
  \sum_{n \in \Z}
  \tbinom{n+m-1}{m}\,
  L_{-n-m} \otimes L_n\,.
\end{equation}
In particular:
\begin{alignat}{2}
  \Omega_0 &= \textstyle\sum_{n} L_{-n} \otimes L_n
  &\qquad& \text{\textup{(}Virasoro Casimir\textup{)}},
  \label{eq:vir-casimir-omega0}
  \\
  \Omega_1 &= \textstyle\sum_{n} n\, L_{-n-1} \otimes L_n\,,
  &&
  \label{eq:vir-omega1}
  \\
  \Omega_2 &= \textstyle\sum_{n}
  \tfrac{n(n+1)}{2}\, L_{-n-2} \otimes L_n\,.
  &&
  \label{eq:vir-omega2}
\end{alignat}
The collision residue recovers the field-valued
form~\eqref{eq:virasoro-r-collision}:
the $z^{-1}$ coefficient is $\Omega_0$, which acts as
$2L$ on any highest-weight module, and the $z^{-3}$ coefficient
is $\Omega_2$, whose action is proportional to $c/2$.
All even-order coefficients $\Omega_{2k}$ for $k \ge 2$ vanish
identically, as does $\Omega_1$, consistent with the
two-term structure~\eqref{eq:virasoro-r-laurent}.
\end{proposition}

\begin{proof}
Write $T(w) = \sum_{n \in \Z} L_n\, w^{-n-2}$ and expand the
collision residue
\[
  r^{\mathrm{Vir}}(z)
  \;=\;
  \frac{c/2}{z^3}
  \;+\;
  \frac{2}{z}\,T
\]
as an operator on $V \otimes V$.  The stress tensor acts on
the first factor by $L_n \otimes 1$; the Casimir-type
coupling to the second factor arises from the OPE origin of
$r(z)$.  Concretely, the OPE
$T(z)\,T(w) \sim \sum_{k=0}^{3} (T_{(k)}T)(w)\,(z{-}w)^{-k-1}$
produces the Laplace kernel
$r^L(z) = \sum_k T_{(k)}T\cdot z^{-k-1}$.
The mode expansion of $r^L(z)$ in
$\mathrm{End}(V)^{\otimes 2}((z^{-1}))$ is obtained by
expanding each $(z{-}w)^{-k-1}$ in the region $|w| < |z|$:
\[
  (z - w)^{-k-1}
  \;=\;
  z^{-k-1}
  \sum_{j \ge 0} \binom{k+j}{j}\,
  \Bigl(\frac{w}{z}\Bigr)^{\!j}.
\]
Extracting the coefficient of $w^{n-1}$ from the first factor
and $w^{-n-1}$ from the second, then shifting by the
$d\log$ pole absorption ($k \mapsto k - 1$), one obtains
the mode coefficients
$\binom{n + m - 1}{m}$ at combined index $m = k - 1$,
yielding~\eqref{eq:vir-r-mode-tensor}.

\medskip
\noindent\textbf{Verification of specialisations.}
At $m = 0$:
$\binom{n-1}{0} = 1$ for all $n$, so
$\Omega_0 = \sum_n L_{-n} \otimes L_n$.
On a highest-weight module of weight $h$, the diagonal
action gives
$\Omega_0\,|h\rangle \otimes |h\rangle
= 2h\,|h\rangle \otimes |h\rangle + \cdots$,
reproducing $r_1 = 2L$.

At $m = 1$:
$\Omega_1 = \sum_n n\, L_{-n-1} \otimes L_n$.
This vanishes as an $r$-matrix coefficient because
the collision residue has no $z^{-2}$ pole: $r_2 = 0$.
Algebraically, $\Omega_1$ contributes to the regular part
of the Laplace kernel (the $\partial T$ term), which is
absorbed by the $d\log$ shift.

At $m = 2$:
$\Omega_2 = \sum_n \frac{n(n+1)}{2}\,L_{-n-2} \otimes L_n$.
On the vacuum $|0\rangle \otimes |0\rangle$, all terms with
$n > 0$ or $n < -2$ annihilate one factor; the surviving
contribution is central and proportional to $c/2$,
recovering $r_3 = c/2$.
\end{proof}

\begin{proposition}[Virasoro twisted coproduct;
\ClaimStatusProvedHere]
\label{prop:virasoro-twisted-coproduct}
The twisted coproduct
$\Delta_z \colon \mathrm{Vir}_c^! \to
\mathrm{Vir}_c^! \otimes_{r(z)} \mathrm{Vir}_c^!$
acts on modes $L_n$, $n \ge 0$, by
\begin{equation}\label{eq:virasoro-coproduct-mode}
  \Delta_z(L_n) \;=\;
  \sum_{k=0}^{n} \binom{n}{k}\, z^k\, L_{n-k} \otimes 1
  \;+\; 1 \otimes L_n\,.
\end{equation}
The $r(z)$-deformed multiplication in the tensor
product acts by the adjoint of $r(z)$ on each factor.
Since $r_3 = c/2$ is central, it commutes with all
modes; the non-trivial contribution comes from
$r_1 = 2L$, which acts via the first $n$-product
$L_{(1)} = 2\,\mathrm{ad}_{L_0}$ of the Virasoro
conformal algebra:
\begin{equation}\label{eq:virasoro-r-adjoint}
  \ad_{r_1/z}(L_n \otimes 1)
  \;=\; \frac{2}{z}\,(L_{(1)} L)_n \otimes 1
  \;=\; -\frac{2n}{z}\,L_n \otimes 1\,,
\end{equation}
using $(L_{(1)} L)_n = (2L)_n = 2L_n$ and the mode
action $\ad_{L_0}(L_n) = -n\,L_n$.
The algebra morphism property
$\Delta_z([L_m, L_n])
= [\Delta_z(L_m),\, \Delta_z(L_n)]_{r(z)}$
holds in $\mathrm{Vir}^! \otimes_{r(z)} \mathrm{Vir}^!$,
with the $z^{-3}$ pole contributing to the central
extension.
\end{proposition}

\begin{proof}
The coproduct~\eqref{eq:virasoro-coproduct-mode} is the
specialisation of the general
formula~\eqref{eq:coproduct-mode} to $a^I_n = L_n$ with
translation~\eqref{eq:translation-modes}.

For the morphism property, expand both sides of
$\Delta_z([L_m, L_n])
= [\Delta_z(L_m), \Delta_z(L_n)]_{r(z)}$
using the Virasoro relation
$[L_m, L_n] = (m{-}n)\,L_{m+n}
+ \tfrac{c}{12}(m^3{-}m)\,\delta_{m+n,0}$.

\medskip
\noindent\textbf{Step 1: The $L_{m+n}$ term.}
Since $\tau_z$ is an algebra automorphism,
$[\tau_z(L_m), \tau_z(L_n)]
= \tau_z([L_m, L_n])$
in the first tensor factor, reproducing
$\tau_z\bigl((m{-}n)\,L_{m+n}\bigr) \otimes 1$.
The second factor gives
$1 \otimes [L_m, L_n]$.  These match the LHS.

\medskip
\noindent\textbf{Step 2: Cross-terms and $r_1$.}
The cross-terms
$\tau_z(L_m) \otimes L_n - \tau_z(L_n) \otimes L_m$
appear on the RHS from expanding the commutator of
$\Delta_z(L_m)$ and $\Delta_z(L_n)$.
In the $r(z)$-deformed tensor product, these are
compensated by $\ad_{r_1/z}$ acting on the
cross-terms.  By~\eqref{eq:virasoro-r-adjoint},
$\ad_{r_1/z}$ contributes
$-2n\,L_n \otimes 1/z$ (respectively
$-2m\,L_m \otimes 1/z$),
and the Vandermonde evaluation of the translation
coefficients
$\sum_{k=0}^{m} \binom{m}{k} z^{k-1}$
produces the correct $(m{-}n)$ factor.

\medskip
\noindent\textbf{Step 3: Central extension and $r_3$.}
The central term
$\tfrac{c}{12}(m^3{-}m)\,\delta_{m+n,0}$
on the LHS is a scalar; it must be reproduced by
the $z^{-3}$ contribution of $r_3 = c/2$ to the
deformed multiplication.  The cross-terms from
Step~2, iterated to cubic order in the translation
coefficients, produce at $z^{-3}$ the sum
\[
  \frac{c}{2z^3}\sum_{k=0}^{m}
  \binom{m}{k}\binom{n}{2-k}\,(-1)^{k+1}
  \;=\;
  \frac{c}{12z^3}(m^3 - m)\,\delta_{m+n,0}\,,
\]
where the identity
$\sum_{k} (-1)^{k+1}\binom{m}{k}\binom{-m}{2-k}
= \tfrac{1}{6}(m^3 - m)$
is verified by direct expansion for
$m + n = 0$, $m \ge 0$.
\end{proof}

\begin{remark}[Non-diagonal mode structure]
\label{rem:virasoro-r-nondiagonal}
The Virasoro $r$-matrix~\eqref{eq:virasoro-r-collision}
differs structurally from the affine
case~\eqref{eq:affine-r-mode} in three ways:
\begin{enumerate}[label=\textup{(\roman*)}]
\item \emph{Pole depth.}
  The affine $r$-matrix has a single pole at $z^{-1}$,
  making the mode coupling diagonal
  ($r_{mn} = k\,\kappa^{IJ}\,\delta_{m+n,0}$).
  The Virasoro has poles at $z^{-1}$ and $z^{-3}$,
  producing couplings at $m + n = 0$ and $m + n = 2$.
  The $z^{-3}$ pole couples mode pairs
  $(L_0, L_2)$, $(L_1, L_1)$, and $(L_2, L_0)$
  through the central charge $c$.
\item \emph{Operator-valued residue.}
  At $z^{-1}$, the affine residue is $k\,\Omega$
  (a c-number Casimir tensor); the Virasoro residue is
  $2L$ (the stress-energy field).  The field $L$ acts
  via $L_0$ on each tensor factor, making the $z^{-1}$
  correction \emph{weight-dependent}: it distinguishes modes
  by their $L_0$-eigenvalue $h - n$ (conformal weight
  minus mode index).
\item \emph{DS origin.}
  The extra pole depth arises from the Drinfeld--Sokolov
  reduction $\widehat{\mathfrak{sl}}_2 \rightsquigarrow
  \mathrm{Vir}_c$.  The affine $r$-matrix
  $r(z) = k\,\Omega/z$ acquires additional poles under DS
  because the Sugawara construction expresses $L$ as a
  quadratic in the currents $J^a$: the quartic Casimir
  $\Omega^{(2)}$ of $\mathfrak{sl}_2$ (which is
  proportional to $(\Omega^{(1)})^2$ by Cayley--Hamilton)
  contributes the $z^{-3}$ pole.
\end{enumerate}
For $W_N$-algebras obtained by DS reduction of
$\widehat{\mathfrak{sl}}_N$, the $r$-matrix acquires poles
at $z^{-1}, z^{-3}, \ldots, z^{-(2N-1)}$, reflecting the
conformal weights $2, 3, \ldots, N$ of the generators.
The mode structure becomes progressively more non-diagonal
with increasing $N$: the shadow depth
$r_{\max} = \infty$ (class~M) is visible in the
pole multiplicity of the $r$-matrix.
\end{remark}

\begin{remark}[Pure Chern--Simons simplification]
\label{rem:pure-cs-simplification}
When $V = 0$ and $k$ is non-degenerate, translations are
$Q$-exact: there exists a derivation $L$ with
$L(B^a(z)^-) = 0$,
$L(c^a(z)^-) = (k^{-1})^{ab}\, B_b(z)^-$,
$[Q, L] = T$,
where $T$ generates holomorphic translation.
The meromorphic dependence of the tensor product
$\otimes_z$ is therefore homotopically trivialised, and one
recovers
$(C_{\mathrm{line}},\, \otimes_z)
\simeq (\cA^!,\, \Delta_z)\text{-}\mathrm{mod}
\simeq \widehat{\fg}_{k - h^\vee}\text{-}\mathrm{mod}$
with the Kac--Moody fusion product.
\end{remark}

\subsubsection{The Chern--Simons twisted coproduct and Maurer--Cartan verification}
\label{subsubsec:cs-twisted-coproduct}

We now specialise the twisted coproduct to Chern--Simons theory
with gauge algebra~$\fg$ and invariant form~$k_{IJ}$, allowing
matter~$V$ (so the effective gauge algebra is
$\mathfrak{h} = \fg \ltimes V$).  The universal coupling
$\mu \in \cA^! \otimes \cA$ and the $r$-matrix
$r(z) \in \cA^! \otimes \cA^!((z^{-1}))$ satisfy independent
Maurer--Cartan equations.  Their compatibility is the
Maurer--Cartan-level proof that $\Delta_z$ is an algebra morphism
into the $r(z)$-deformed tensor product.

\begin{proposition}[Chern--Simons twisted coproduct in mode form;
\ClaimStatusProvedHere]
\label{prop:cs-twisted-coproduct}
Let\/ $\fg$ be a finite-dimensional Lie algebra with
non-degenerate invariant form\/ $k_{IJ}$, and let\/
$\cA^! = \cA^!_{\mathrm{line}}$ be the open-colour Koszul dual
of the HT Chern--Simons theory, generated by\/ $B^I_n$
\textup{(}even\textup{)} and\/ $c^I_n$
\textup{(}odd\textup{)}, $n \ge 0$, $1 \le I \le \dim \mathfrak{h}$,
with bracket, differential, and $r$-matrix as in
Proposition~\textup{\ref{prop:coproduct-morphism}}.  The twisted
coproduct\/
$\Delta_z \colon \cA^! \to \cA^! \otimes_{r(z)} \cA^!$
acts on generators by
\begin{align}
  \Delta_z(B^I_n)
  &\;=\;
  \sum_{r=0}^{n} \binom{n}{r}\, z^r\, B^I_{n-r}
  \;\otimes\; 1
  \;+\;
  1 \;\otimes\; B^I_n,
  \label{eq:cs-coproduct-B}
  \\[4pt]
  \Delta_z(c^I_n)
  &\;=\;
  \sum_{r=0}^{n} \binom{n}{r}\, z^r\, c^I_{n-r}
  \;\otimes\; 1
  \;+\;
  1 \;\otimes\; c^I_n.
  \label{eq:cs-coproduct-c}
\end{align}
In generating-function form,
\begin{equation}\label{eq:cs-coproduct-gf}
  \Delta_z\bigl(B^I(w)^-\bigr)
  \;=\; B^I(w+z)^- \otimes 1 \;+\; 1 \otimes B^I(w)^-,
  \qquad
  \Delta_z\bigl(c^I(w)^-\bigr)
  \;=\; c^I(w+z)^- \otimes 1 \;+\; 1 \otimes c^I(w)^-.
\end{equation}
The target of\/ $\Delta_z$ is the\/ $r(z)$-deformed tensor
product\/ $\cA^! \otimes_{r(z)} \cA^!$, whose multiplication is
$\mu_{\otimes} = \mu_1 \otimes 1 + 1 \otimes \mu_2
+ \ad_{r(z)}$; the $r$-matrix correction\/
$r(z) = k\,\Omega/z$
\textup{(}Proposition~\textup{\ref{prop:affine-r-mode}}\textup{)}
deforms the cross-terms exactly so that\/ $\Delta_z$ is an algebra
morphism.
\end{proposition}

\begin{proof}
By Proposition~\ref{prop:twisted-coproduct-modes}, the formulas
\eqref{eq:cs-coproduct-B}--\eqref{eq:cs-coproduct-c} are the
specialisation of the general coproduct
$\Delta_z(a^I_n) = \tau_z(a^I_n) \otimes 1 + 1 \otimes a^I_n$
to the generators $(a^I) = (B^I, c^I)$.  The
generating-function form~\eqref{eq:cs-coproduct-gf} restates
$\tau_z(B^I(w)^-) = B^I(w+z)^-$, which is the definition of the
translation automorphism.  That $\Delta_z$ lands in the
$r(z)$-deformed tensor product is the content of
Proposition~\ref{prop:coproduct-morphism}:
the cross-terms $\tau_z(B^I_m) \otimes B^J_n
- \tau_z(B^J_n) \otimes B^I_m$ arising in the commutator are
absorbed by $\ad_{r(z)}$ via $[r(z),\, B^J_n \otimes 1]
= k\,f^{KJ}{}_M\,\kappa^{KL}\,B^M_n \otimes t_L/z$.
\end{proof}

\begin{proposition}[Maurer--Cartan verification of the coproduct;
\ClaimStatusProvedHere]
\label{prop:mc-coproduct-verification}
Define the universal coupling element
\begin{equation}\label{eq:universal-coupling}
  \mu
  \;=\;
  \oint_u \bigl(
    c^I(u)^-\, B_I(u)^+
    \;+\;
    B^I(u)^-\, c_I(u)^+
  \bigr),
\end{equation}
where\/ $B_I(u)^+ = k_{IJ}\,B^J(u)^+$ and\/
$c_I(u)^+ = k_{IJ}\,c^J(u)^+$ are the positive-mode fields of
the closed-colour algebra\/ $\cA$, and the integral extracts the
residue at\/ $u = 0$.  Let\/
$\mu(z) = (\tau_z \otimes \id)(\mu)$ denote the translated
coupling and\/ $\mu'(0) = (1 \otimes \mu)$ the second-copy
coupling.  Then:
\begin{enumerate}[label=\textup{(\roman*)}]
\item $\mu$ satisfies the Maurer--Cartan equation\/
  $Q\mu + \tfrac{1}{2}[\mu, \mu] = 0$
  in\/ $\cA^! \otimes \cA$.
\item $r(z)$ satisfies the classical Yang--Baxter equation\/
  $[r_{12}(z_{12}),\, r_{13}(z_{13})]
  + [r_{12}(z_{12}),\, r_{23}(z_{23})]
  + [r_{13}(z_{13}),\, r_{23}(z_{23})] = 0$.
\item The combined element\/
  $\mu(z) + \mu'(0) + r(z)$
  is Maurer--Cartan in\/
  $\cA^! \otimes_{r(z)} \cA^! \otimes \cA$:
\begin{equation}\label{eq:combined-mc}
  [\mu(z),\, \mu'(0)]
  \;+\; [r(z),\, \mu'(0)]
  \;+\; [r(z),\, \mu(z)]
  \;=\; 0.
\end{equation}
\end{enumerate}
Equation~\eqref{eq:combined-mc} is the
Maurer--Cartan-level statement that\/ $\Delta_z$ is an\/
$A_\infty$ algebra morphism and that the OPE is associative.
\end{proposition}

\begin{proof}
\textbf{(i).}
The coupling $\mu$ is the residue pairing between dual generators:
$\mu = \sum_{n \ge 0}(c^I_n \otimes B_{I,n} + B^I_n \otimes c_{I,n})$,
where $(B_{I,n}, c_{I,n})$ are the positive modes of $\cA$.
The Maurer--Cartan equation $Q\mu + \tfrac{1}{2}[\mu, \mu] = 0$ is
the statement that $\mu$ is a twisting morphism from the coalgebra
$\cA^!$ to the algebra $\cA$.  By the bar-cobar adjunction
(Volume~I, Theorem~A), twisting morphisms biject with
dg-algebra maps $\Omega(B(\cA)) \to \cA$, and the universal
twisting morphism is Maurer--Cartan by construction.

\smallskip\noindent
\textbf{(ii).}
This is the content of
Proposition~\ref{prop:field-theory-r}, Step~4: the CYBE follows
from the order-$\hbar^2$ expansion of the quantum Yang--Baxter
equation (Theorem~\ref{thm:YBE}).

\smallskip\noindent
\textbf{(iii).}
The three brackets in~\eqref{eq:combined-mc} arise from expanding
the full Maurer--Cartan equation for the element
$\Phi = \mu(z) + \mu'(0) + r(z)$ in the convolution dg Lie
algebra $\cA^! \otimes_{r(z)} \cA^! \otimes \cA$.
Since $\mu(z)$ and $\mu'(0)$ live in separate tensor slots (the
first and second copy of $\cA^!$), their mutual bracket reduces to
the cross-term of the coproduct.  Since $r(z) \in \cA^! \otimes \cA^!$
and $\mu \in \cA^! \otimes \cA$, the brackets
$[r(z), \mu'(0)]$ and $[r(z), \mu(z)]$ pair one slot of $r(z)$
against one slot of $\mu$ via the $\fg$-action.

Computing explicitly: the bracket $[\mu(z), \mu'(0)]$ produces
cross-channel OPE singularities: a field $c^I(u+z)^-$ from
$\mu(z)$ contracts with $B_I(v)^+$ from $\mu'(0)$, yielding a
pole at $u + z = v$.  The bracket $[r(z), \mu'(0)]$ produces a
$1/z$ pole from the contraction of $r(z) = k\,\Omega/z$ with the
$\fg$-valued part of $\mu'(0)$.  The bracket $[r(z), \mu(z)]$
produces the corresponding correction from the first copy.
The sum vanishes by the Jacobi identity for $\fg$ applied to the
structure constants: the coefficient of each pole in $z$ and $u$
is proportional to
$f^{IJ}{}_K f^{KL}{}_M + f^{JL}{}_K f^{IK}{}_M
+ f^{LI}{}_K f^{JK}{}_M = 0$.

The vanishing of~\eqref{eq:combined-mc} states that $\Phi$ is
Maurer--Cartan, i.e.\ the two-copy coupling
$\mu(z) + \mu'(0)$ is compatible with the $r(z)$-deformation.
In the language of $A_\infty$ morphisms, this is the
equation $\Delta_z \circ m_k = m_k^{\otimes} \circ \Delta_z^{\otimes k}$
at the universal level: the coproduct intertwines the $A_\infty$
operations of $\cA^!$ with those of
$\cA^! \otimes_{r(z)} \cA^!$.
\end{proof}

\begin{remark}[Pure Chern--Simons and current modes]
\label{rem:cs-current-coproduct}
For pure Chern--Simons ($V = 0$, $\mathfrak{h} = \fg$), the
$Q$-cohomology of $\cA^!$ is generated by the current modes
$J^a_n := (k^{-1})^{ab}\,B_{b,n}$, $n \ge 0$,
$1 \le a \le \dim\fg$, and the ghost generators $c^a_n$ become
$Q$-exact.  The coproduct~\eqref{eq:cs-coproduct-B} descends to
\begin{equation}\label{eq:current-coproduct}
  \Delta_z(J^a_n)
  \;=\;
  \sum_{r=0}^{n} \binom{n}{r}\, z^r\, J^a_{n-r}
  \;\otimes\; 1
  \;+\;
  1 \;\otimes\; J^a_n,
\end{equation}
with the $r(z)$-correction now acting through
$[r(z),\, 1 \otimes J^a_n]
= k\,\kappa^{bc}\,f_{bc}{}^a\,J^a_n \otimes t_c/z$.
Since translations are $Q$-exact in the topological theory
(Remark~\ref{rem:pure-cs-simplification}), the meromorphic
dependence on~$z$ is homotopically trivial, and
$\Delta_z$ reduces on $Q$-cohomology to the
standard Yangian coproduct $\Delta_z(J^a_n)
= \tau_z(J^a_n) \otimes 1 + 1 \otimes J^a_n$
with target the \emph{undeformed} tensor product of the
Kac--Moody algebra $\widehat{\fg}_{k - h^\vee}$.
\end{remark}

\subsection{dg-Shifted Yangians}
\label{subsec:dg-yangian-operadic}

\begin{definition}[dg-Shifted Yangian]
\label{def:dg_Yangian}
A \emph{dg-shifted Yangian} is an $A_\infty$ algebra $\mathcal{Y}$ equipped with:
\begin{enumerate}[label=(\roman*)]
\item A translation automorphism $\tau_z : \mathcal{Y} \to \mathcal{Y}[[z]]$, generated by a derivation $T$ of degree $(0, \text{odd}, 1)$:
\begin{equation}
\tau_z = \exp(z T) = \sum_{n \geq 0} \frac{z^n}{n!} T^n;
\end{equation}

\item A spectral $R$-matrix $r(z) \in \mathcal{Y} \widehat{\otimes} \mathcal{Y}[[z^{-1}, z]]$ of degree $(1, \text{odd}, 0)$;

\item A twisted coproduct $\Delta_z : \mathcal{Y} \to \mathcal{Y} \otimes_{r(z)} \mathcal{Y}[[z^{-1}, z]]$ (an $A_\infty$ algebra morphism);
\item A counit $\varepsilon : \mathcal{Y} \to \C$ (also an $A_\infty$ morphism);
\end{enumerate}
satisfying:
\begin{itemize}
\item Counit axioms: $\varepsilon(1) = 1$, $(\varepsilon \otimes id) \circ \Delta_z = id$, $(id \otimes \varepsilon) \circ \Delta_z \simeq \tau_z$;
\item Translation invariance: $\varepsilon(T(a)) = 0$, $\varepsilon(\tau_z(a)) = \varepsilon(a)$ for all $a \in \mathcal{Y}$;
\item Weak co-commutativity: $(\tau_w \otimes \tau_w) r(z) = r(z)$, $r(-z) = -r_{21}(z)$;
\item Co-associativity (quantum $A_\infty$ Yang--Baxter):
\begin{equation}
r_{23}(z) + (id \otimes \Delta_z)(r(z+w)) = r_{12}(w) + (\Delta_w \otimes id)(r(z)).
\end{equation}
\end{itemize}
\end{definition}

\begin{remark}[Antipode and Hopf structure]
\label{rem:dg-yangian-antipode}
Definition~\ref{def:dg_Yangian} produces a dg-shifted \emph{bialgebra} (coproduct $\Delta_z$, counit $\varepsilon$), not a Hopf algebra: no antipode is defined. The Drinfeld Yangian $Y(\fg)$ has an antipode $S$; in the $A_\infty$ setting, the existence of a compatible antipode is an open question. For the representation-theoretic applications in this monograph (line-operator modules, spectral braiding), the bialgebra structure suffices.
\end{remark}

\begin{theorem}[Affine Open-Colour Dual is a dg-Shifted Yangian;
\ClaimStatusProvedHere]
\label{thm:Koszul_dual_Yangian}
(Dimofte--Niu--Py \cite{DNP25}, Theorem 5.5) For the standard affine
HT gauge realization with closed colour $V_k(\fg)$, satisfying the
hypotheses of Theorem~\ref{thm:physics-bridge}, the open-colour
Koszul dual algebra $\mathcal{A}^!_{\mathrm{line}}$ is a
dg-shifted Yangian in the sense of Definition~\ref{def:dg_Yangian}.
\end{theorem}

\begin{proof}
We verify each piece of structure in
Definition~\ref{def:dg_Yangian} for this affine HT gauge
realization.

\medskip
\noindent\textbf{Step 1: Translation operator.}
The holomorphic translation $z \mapsto z + w$ acts on bulk local operators of the 3d HT theory.  Under the bar--cobar adjunction (Theorem~\ref{thm:bar-cobar-adjunction}), this induces a derivation $T$ of degree $(0,\text{odd},1)$ on the open-colour Koszul dual $\A^!_{\mathrm{line}}$: if $\phi(z)$ is a boundary-localized field, then $T\phi = \partial_z \phi$.  The automorphism $\tau_w = \exp(wT) = \sum_{n \geq 0} w^n T^n / n!$ converges in $\A^!_{\mathrm{line}}[[w]]$ because $T$ raises filtration degree (the holomorphic weight filtration of Proposition~\ref{prop:gr-chiral}).

\medskip
\noindent\textbf{Step 2: Spectral $R$-matrix.}
The OPE of two parallel line operators $L_1(z_1)$, $L_2(z_2)$ is governed by configuration integrals over $\FM_2(\C) \times E_1(2)$.  Setting $z = z_1 - z_2$, the quantum corrections to the free tensor product define $r(z) \in \A^!_{\mathrm{line}} \widehat{\otimes}\, \A^!_{\mathrm{line}}((z^{-1}))$ as in Theorem~\ref{thm:spectral_R_YBE}.  The degree is $(1, \text{odd}, 0)$: degree 1 from the desuspension in the bar construction, odd under the $\Z/2$-grading from the fermion number, and weight 0 because the leading pole $1/z$ carries no holomorphic weight shift.

\medskip
\noindent\textbf{Step 3: Coproduct from the monoidal structure on $\mathcal{C}_{\mathrm{line}}$.}
Placing two parallel lines at $z_1, z_2 \in \C$ and taking the OPE defines the coproduct.  Concretely, for $x \in \A^!_{\mathrm{line}}$, the element $\Delta_z(x)$ encodes how the action of $x$ on a single line decomposes when the line is resolved into a pair separated by spectral parameter $z$:
\[
\Delta_z \colon \A^!_{\mathrm{line}} \longrightarrow \A^!_{\mathrm{line}} \otimes_{r(z)} \A^!_{\mathrm{line}}[[z^{-1}, z]],
\]
where the twisted tensor product $\otimes_{r(z)}$ has differential $d_{\otimes} + [r(z), -]$.  This is an $A_\infty$ algebra morphism: the compatibility with higher operations $m_k$ follows from the $W(\SCchtop)$-module structure on boundary factorization modules (Section~\ref{subsec:boundary-module}), which ensures the OPE respects the full homotopy-coherent algebra structure.

The $R$-matrix intertwines the coproduct and its opposite:
\[
R(z)\,\Delta_z(x) = \Delta_z^{\mathrm{op}}(x)\, R(z) \qquad \text{for all } x \in \A^!_{\mathrm{line}}.
\]
This is the defining property of a quasi-triangular structure.  It follows from the definition of $R(z)$ as the braiding isomorphism (Definition~\ref{def:spectral-braiding}): applying $R(z)$ swaps the two tensor factors, converting $\Delta_z$ (first factor at $z_1$, second at $z_2$) to $\Delta_z^{\mathrm{op}}$ (factors swapped).

\medskip
\noindent\textbf{Step 4: Counit.}
Define $\varepsilon \colon \A^!_{\mathrm{line}} \to \C$ by the augmentation: removing a line from a configuration corresponds to the counit of the $E_1$-coalgebra structure.  Explicitly, $\varepsilon$ evaluates a line operator on the vacuum module.  Then:
\begin{itemize}
\item $\varepsilon(1) = 1$ (the identity line evaluates to the scalar identity);
\item $(\varepsilon \otimes \id) \circ \Delta_z = \id$ (removing the first line recovers the second);
\item $(\id \otimes \varepsilon) \circ \Delta_z \simeq \tau_z$ (removing the second line at position $z_2$ and keeping the first at $z_1$ shifts by $z = z_1 - z_2$ via holomorphic translation).
\end{itemize}

\medskip
\noindent\textbf{Step 5: Translation invariance.}
Simultaneous holomorphic translation $z_i \mapsto z_i + w$ for all lines preserves the spectral differences $z_{ij}$.  Therefore:
\begin{itemize}
\item $\varepsilon(T(a)) = \varepsilon(\partial_z a) = 0$ (the augmentation kills derivatives, since the vacuum is translation-invariant);
\item $(\tau_w \otimes \tau_w)\, r(z) = r(z)$ (the $R$-matrix depends only on $z = z_1 - z_2$, not on the individual positions).
\end{itemize}

\medskip
\noindent\textbf{Step 6: Co-associativity (quantum $A_\infty$ Yang--Baxter).}
The identity
\[
r_{23}(z) + (\id \otimes \Delta_z)(r(z+w)) = r_{12}(w) + (\Delta_w \otimes \id)(r(z))
\]
expresses the consistency of the three-line OPE: resolving the pair $(L_2, L_3)$ first and then braiding $L_1$ past the pair (LHS) must agree with braiding $L_1$ past $L_2$ first and then resolving (RHS).  This is a rewriting of the Yang--Baxter equation (Theorem~\ref{thm:YBE}) in coproduct language, obtained by substituting $R(z) = \exp(r(z))$ and expanding the three-fold braiding identity to leading nontrivial order.

\medskip
\noindent\textbf{Step 7: Yangian structure.}
The spectral parameter $z$ makes this a \emph{Yangian} rather than a mere quasi-triangular Hopf algebra: $\Delta_z$, $r(z)$, and $\tau_z$ assemble into a deformation of $U(\mathfrak{g}[z])$ (the universal enveloping algebra of the loop algebra).  In the associated graded with respect to the holomorphic filtration, $r(z) \sim r_0/z$ (simple pole) and $\Delta_z$ reduces to the standard Yangian coproduct $\Delta_z(x) = x \otimes 1 + 1 \otimes x + O(z^{-1})$.  This matches Definition~\ref{def:dg_Yangian}.
\end{proof}

\begin{theorem}[Recognition on the chirally Koszul locus;
\ClaimStatusProvedHere]
\label{thm:yangian-recognition}
\index{recognition theorem!Swiss-cheese to Yangian|textbf}
\index{dg-shifted Yangian!recognition|textbf}
Let\/ $\cA$ be a chirally Koszul logarithmic
$\SCchtop$-algebra, and let\/ $\cA^!_{\mathrm{line}}$ denote
its open-colour Koszul dual.  Then
$\cA^!_{\mathrm{line}}$ is a dg-shifted Yangian in the sense
of Definition~\textup{\ref{def:dg_Yangian}}.
\end{theorem}

\begin{proof}
For the forward direction, Theorem~\ref{thm:dual-sc-algebra}
supplies the open-colour Yangian product, the closed-colour
Sklyanin bracket, and the holomorphic translation automorphism
on~$\cA^!_{\mathrm{line}}$ (the last being
axiom~(i) of Definition~\ref{def:dg_Yangian}).
Theorem~\ref{thm:spectral-parameter-from-closed-color}
identifies the spectral parameter with the closed colour,
Theorem~\ref{thm:YBE} gives the quantum $A_\infty$
Yang--Baxter relation, and on the chirally Koszul locus
Theorem~\ref{thm:lines_as_modules} models the line category
by $\cA^!_{\mathrm{line}}$-modules, so the monoidal
line category
provides the twisted coproduct and counit.  These data satisfy
Definition~\ref{def:dg_Yangian}, hence
$\cA^!_{\mathrm{line}}$ is a
dg-shifted Yangian.
\end{proof}

\begin{remark}[Scope of the recognition theorem]
\label{rem:yangian-recognition-scope}
The theorem used on the live surface is the forward direction:
on the chirally Koszul locus, the open-colour Koszul dual
carries the dg-shifted Yangian package.  A converse
reconstruction from an abstract dg-shifted Yangian back to a
logarithmic $\SCchtop$-algebra would require an independent
recovery of the full mixed-colour operadic structure; that
stronger converse is not used elsewhere in the manuscript.
\end{remark}

\subsection{Connection to Latyntsev's factorisation quantum groups}

The spectral $R$-matrix and the meromorphic tensor product are an instance of Latyntsev's \emph{factorisation quantum groups} \cite{Lat23}.

\begin{definition}[Factorisation Quantum Group (Informal)]
\label{def:factorisation_QG}
A \emph{factorisation quantum group} is a categorical structure encoding:
\begin{itemize}
\item A monoidal category $\mathcal{C}$ of representations;
\item A meromorphic tensor product $\otimes_z : \mathcal{C} \times \mathcal{C} \to \mathcal{C}[[z^{-1}, z]]$, parametrized by a spectral parameter $z \in \C$;
\item An $R$-matrix $R(z) : V \otimes_z W \to W \otimes_{-z} V$ satisfying the quantum Yang--Baxter equation.
\end{itemize}
\end{definition}

\begin{theorem}[Line Operators Form a Factorisation Quantum Group;
\ClaimStatusProvedHere]
\label{thm:lines_factorisationQG}
Let $A$ be a logarithmic $\SCchtop$-algebra, and suppose its
boundary line operators carry the meromorphic OPE and spectral
braiding of Definition~\ref{def:spectral-braiding}; in particular,
this applies to physical realizations satisfying
Theorem~\ref{thm:physics-bridge}.  Then the resulting category
$\mathcal{C}_{\text{line}}$ forms a factorisation quantum group in
the sense of Latyntsev \cite{Lat23}.
\end{theorem}

\begin{proof}
We verify the three axioms of a factorisation quantum group (Definition~\ref{def:factorisation_QG}).

\medskip
\noindent\textbf{(i) Monoidal structure.}
The category $\mathcal{C}_{\mathrm{line}}$ has objects given by
boundary line operators and morphisms given by boundary-localized
local operators (intertwiners).  The monoidal product is
concatenation in the $E_1$-direction along the boundary $t = 0$:
given lines $L_1, L_2$ placed at distinct points in $\C$, the
product $L_1 \otimes L_2$ is defined by the $W(\SCchtop)$-module
structure (Section~\ref{subsec:boundary-module}).  Associativity and
unit constraints follow from the $E_1$-coherence encoded in the
$W$-resolution.  On the chirally Koszul locus,
Theorem~\ref{thm:lines_as_modules} further identifies
$\mathcal{C}_{\mathrm{line}}$ with the category of
$\A^!_{\mathrm{line}}$-modules, but that extra comparison is not needed for the
present factorisation-quantum-group axioms.

\medskip
\noindent\textbf{(ii) Meromorphic tensor product.}
For each $z \in \C^\times$, the OPE of line operators at separation $z$ defines a meromorphic family of tensor functors
\[
\otimes_z \colon \mathcal{C}_{\mathrm{line}} \times \mathcal{C}_{\mathrm{line}} \longrightarrow \mathcal{C}_{\mathrm{line}}[[z^{-1}, z]].
\]
Concretely, $L_1 \otimes_z L_2$ is the line operator obtained by
placing $L_1$ at position $z$ and $L_2$ at position $0$ and taking
the bulk-to-boundary OPE.  The OPE coefficients $C_k(z)$ are
meromorphic in $z$ with poles only at $z = 0$ (collision
singularities), as supplied in physical realizations by
Theorem~\ref{thm:physics-bridge}.  The family $\otimes_z$ varies
holomorphically in $z$ on $\C^\times$: by
Definition~\ref{def:log-SC-algebra} and
Theorem~\ref{thm:FM-calculus}, the FM-compactified configuration
integrals that define the OPE are smooth functions of the spectral
parameter away from the collision divisor.

\medskip
\noindent\textbf{(iii) Braiding by $R(z)$.}
The spectral braiding $R(z) \colon L_1 \otimes_z L_2 \xrightarrow{\sim} L_2 \otimes_{-z} L_1$ (Definition~\ref{def:spectral-braiding}) is an isomorphism between $\otimes_z$ and the opposite product $\otimes_{-z}^{\mathrm{op}}$.  By Theorem~\ref{thm:YBE}, $R(z)$ satisfies the quantum Yang--Baxter equation, i.e.\ the braid relation for a braided monoidal category with spectral parameter.

\medskip
\noindent\textbf{(iv) Factorisation property.}
The distinguishing feature of a \emph{factorisation} quantum group (as opposed to an ordinary quantum group) is that the braiding depends holomorphically on $z$, upgrading the braided monoidal category to a \emph{factorisation category} on $\C$.  In Latyntsev's framework \cite{Lat23}, this means: the assignment $z \mapsto (\mathcal{C}_{\mathrm{line}}, \otimes_z, R(z))$ defines a constructible cosheaf on the Ran space $\Ran(\C)$.  For our category, this factorisation structure is inherited from the prefactorization algebra $\Obs^{\partial}$ on $\C \times \R_{\geq 0}$: restricting $\Obs^{\partial}$ to disks at different points of $\C$ and taking the $E_1$-coinvariants in the $\R$-direction yields fibers of the cosheaf, and the factorisation isomorphisms come from the prefactorization algebra structure maps for disjoint disks.

The holomorphic variation of the braiding in $z$ is therefore a direct consequence of the holomorphic-topological factorisation algebra structure of the 3d HT theory, restricted to the boundary.
\end{proof}

\subsection{Example: abelian CS and Heisenberg Yangian}

\begin{example}[Heisenberg Yangian from $U(1)$ CS]
\label{ex:Heisenberg_Yangian}
For abelian $U(1)$ Chern--Simons at level~$k$
\textup{(}Example~\textup{\ref{ex:abelian_CS_complete}}\textup{)},
the bulk algebra is the affine Kac--Moody
$\widehat{\mathfrak{u}(1)}_k$, and the open-colour Koszul dual is
the Heisenberg Yangian $Y(\mathfrak{u}(1))$.  The spectral
$R$-matrix is
\begin{equation}
r(z) = k \, \frac{J \otimes J}{z} + \text{(higher terms)},
\end{equation}
satisfying the classical Yang--Baxter equation.  For non-abelian $G$,
the open-colour dual is no longer abelian: in the standard affine HT
gauge realization it is the dg-shifted Yangian $\Ydg(\mathfrak{g})$.
\end{example}

%%%%%%%%%%%%%%%%%%%%%%%%%%%%%%%%%%%%%%%%%%%%%%%%%%%%%%%%%%%%%%%%%%%%%%%%%%%%%
% THE E_1-CHIRAL STRUCTURE OF THE SHIFTED YANGIAN
%%%%%%%%%%%%%%%%%%%%%%%%%%%%%%%%%%%%%%%%%%%%%%%%%%%%%%%%%%%%%%%%%%%%%%%%%%%%%

\subsection{The $E_1$-chiral structure of the shifted Yangian}
\label{subsec:E1-chiral-yangian}
\index{dg-shifted Yangian!$E_1$-chiral structure|textbf}
\index{Yangian!spectral OPE}
\index{Yangian!PVA descent}
\index{Sklyanin bracket|textbf}

On the chirally Koszul locus, the open-colour Koszul dual
$\cA^!_{\mathrm{line}}$ is a
dg-shifted Yangian
\textup{(}Theorem~\textup{\ref{thm:yangian-recognition}}\textup{)},
and line operators are $\cA^!_{\mathrm{line}}$-modules
(Theorem~\ref{thm:lines_as_modules}).
But $\cA^!_{\mathrm{line}}$ carries a richer structure: it is an $E_1$-chiral
algebra on $\C_u \times \R_t$, with the $E_1$ direction from
the deconcatenation coproduct on $\R_t$ and the chiral direction
from the spectral OPE on $\C_u$.  We develop this structure
from the RTT relation, working from explicit formulas; the
definitive operadic framework is the master two-color Koszul
duality theorem (Theorem~\ref{thm:two-color-master}).

\subsubsection*{The FRT spectral OPE}

The starting point is the FRT commutation relation, derived
from the RTT relation $R_{12}(u{-}v)\,T_1(u)\,T_2(v) =
T_2(v)\,T_1(u)\,R_{12}(u{-}v)$
(equation~\eqref{eq:rtt-from-bar}).  For $\fg$ simple with
rational $R$-matrix $R(u) = 1 + \hbar\,P/u$ ($P$ the
permutation on $V \otimes V$, $V$ the defining representation),
the component form is
\begin{equation}\label{eq:FRT-components}
\boxed{
[t_{ij}(u),\, t_{kl}(v)]
\;=\;
\frac{\hbar}{u - v}\bigl(
  t_{kj}(u)\,t_{il}(v)
  \;-\;
  t_{kj}(v)\,t_{il}(u)
\bigr),
}
\end{equation}
where $T(u) = \sum_{I,J} E_{IJ} \otimes t_{IJ}(u) \in
\operatorname{End}(V) \otimes Y_\hbar(\fg)[\![u^{-1}]\!]$.
This is the \emph{spectral OPE}: it has the structure of a
chiral bracket with singularity at $u = v$.  The right-hand
side is a \emph{difference} of products evaluated at $u$ and
$v$, and vanishes at $u = v$, so the pole $1/(u-v)$ is
cancelled, making the bracket well-defined as a formal
distribution.

\subsubsection*{Calibration: the Heisenberg case}

\begin{computation}[Heisenberg Koszul dual as $E_1$-chiral
algebra; \ClaimStatusProvedHere]
\label{comp:heisenberg-E1-chiral}
For $\cA = \widehat{\mathfrak{u}(1)}_k$, the open-colour Koszul
dual is
$\cA^!_{\mathrm{line}} = Y(\mathfrak{u}(1))
\simeq \widehat{\mathfrak{u}(1)}_{-k}$.
There is a single FRT generator $t(u)$,
and the $R$-matrix is $R(u) = 1 + k/u$ (scalar, since
$\dim V = 1$).  The FRT relation~\eqref{eq:FRT-components}
gives
\[
[t(u),\, t(v)]
\;=\;
\frac{k}{u - v}\bigl(
  t(u)\,t(v) - t(v)\,t(u)
\bigr)
\;=\;
\frac{k}{u - v}\,[t(u), t(v)].
\]
This forces $[t(u), t(v)] = 0$ for $k \ne 0$
\textup(the zero mode is abelian\textup).
Passing to the current $J(u) = \partial_u \log t(u)$: the
relation becomes
\begin{equation}\label{eq:heisenberg-spectral-OPE}
[J_m,\, J_n] \;=\; -k\,m\,\delta_{m+n,0},
\end{equation}
and the spectral OPE is
$J(u)\,J(v) \sim -k/(u-v)^2$.
This is the Heisenberg algebra at level $-k$: the Koszul dual
flips the sign of the level.

The $E_1$-chiral structure is manifest: the spectral OPE has a
double pole at $u = v$ (a ``chiral'' singularity in the spectral
variable), and the algebra on $\R_t$ is the ordinary
Heisenberg, associative with no higher homotopies.  The
$A_\infty$ operations $m_k = 0$ for $k \ge 3$: the Heisenberg
Yangian is \emph{formal}.  This is the abelian calibration.
\end{computation}

\subsubsection*{The $\mathfrak{sl}_2$ computation}

For $\fg = \mathfrak{sl}_2$, the transfer matrix is
$T(u) = \bigl(\begin{smallmatrix}
t_{11}(u) & t_{12}(u) \\ t_{21}(u) & t_{22}(u)
\end{smallmatrix}\bigr)$
with $R(u) = 1 + \hbar P/u$,
$P = \sum_{i,j} E_{ij} \otimes E_{ji}$.
Equation~\eqref{eq:FRT-components} gives $16$ relations, of
which the independent ones are \textup(writing
$w = u - v$\textup):
\begin{align}
[t_{11}(u),\, t_{21}(v)]
&\;=\;
\frac{\hbar}{w}\bigl(t_{21}(u)\,t_{11}(v)
  - t_{21}(v)\,t_{11}(u)\bigr),
\label{eq:sl2-RTT-1}
\\[3pt]
[t_{21}(u),\, t_{12}(v)]
&\;=\;
\frac{\hbar}{w}\bigl(t_{11}(u)\,t_{22}(v)
  - t_{11}(v)\,t_{22}(u)\bigr),
\label{eq:sl2-RTT-2}
\\[3pt]
[t_{21}(u),\, t_{21}(v)]
&\;=\;
\frac{\hbar}{w}\bigl(t_{21}(u)\,t_{21}(v)
  - t_{21}(v)\,t_{21}(u)\bigr)
\;=\; 0.
\label{eq:sl2-RTT-3}
\end{align}
Relation~\eqref{eq:sl2-RTT-2} is the crucial one: it produces
the $[E, F] \sim H$ commutation relation of the Yangian.

In the Gauss decomposition
$T(u) = \bigl(\begin{smallmatrix} 1 & 0 \\ E(u) & 1
\end{smallmatrix}\bigr)
\bigl(\begin{smallmatrix} k_1(u) & 0 \\ 0 & k_2(u)
\end{smallmatrix}\bigr)
\bigl(\begin{smallmatrix} 1 & F(u) \\ 0 & 1
\end{smallmatrix}\bigr)$,
the Drinfeld generators $(E(u), F(u), H(u) = k_1(u)/k_2(u))$
satisfy
\begin{align}
(u - v - \hbar)\,E(u)\,E(v)
&\;=\;
(u - v + \hbar)\,E(v)\,E(u),
\label{eq:drinfeld-EE}
\\[3pt]
[E(u),\, F(v)]
&\;=\;
\frac{\hbar}{u - v}
\bigl(k_1(u)\,k_2(v) - k_1(v)\,k_2(u)\bigr),
\label{eq:drinfeld-EF}
\end{align}
derived from~\eqref{eq:sl2-RTT-1}--\eqref{eq:sl2-RTT-2} by
substitution.  These are the \emph{exact} spectral OPE
relations, not truncated at leading order.

\begin{remark}[What the spectral OPE is and is not]
\label{rem:spectral-OPE-nature}
Equations~\eqref{eq:FRT-components}
and~\eqref{eq:drinfeld-EE}--\eqref{eq:drinfeld-EF} are
\emph{not} $\lambda$-brackets in the PVA sense.  They are
commutation relations between generating functions, with
singularity at $u = v$.  The term ``spectral OPE'' is used
by analogy: the pole structure at $u = v$ parallels the pole
structure of a chiral algebra OPE at $z = w$, and the
residues encode the ``$E_1$-chiral product'' of line operators
at coincident spectral position.  But the generating functions
$E(u)$, $F(u)$, $H(u)$ are \emph{not} conformal fields: they
have no conformal weight, no state-field correspondence, and
no normally-ordered product.  The $E_1$-chiral structure on the
Yangian is a \emph{filtered} chiral algebra in the spectral
variable, not a vertex algebra.
\end{remark}

\subsubsection*{The classical limit: Sklyanin--Drinfeld Poisson
bracket}

\begin{proposition}[Sklyanin bracket from classical RTT;
\ClaimStatusProvedHere]
\label{prop:sklyanin-from-RTT}
\index{Sklyanin bracket|textbf}
The classical limit $\hbar \to 0$ of the FRT
relation~\eqref{eq:FRT-components}, with the replacement
$[\cdot, \cdot]/\hbar \to \{\cdot, \cdot\}$, gives
\begin{equation}\label{eq:sklyanin-bracket}
\{t_{ij}(u),\, t_{kl}(v)\}
\;=\;
\frac{1}{u - v}\bigl(
  t_{kj}(u)\,t_{il}(v) - t_{kj}(v)\,t_{il}(u)
\bigr).
\end{equation}
This is the \emph{Sklyanin--Drinfeld Poisson bracket} on the
classical Yangian $Y_{\mathrm{cl}}(\fg) \simeq
\mathcal{O}(G[\![u^{-1}]\!])$, the coordinate ring of the
formal arc group.
\end{proposition}

\begin{proof}
Divide~\eqref{eq:FRT-components} by $\hbar$ and take
$\hbar \to 0$.  The left side
$[t_{ij}(u), t_{kl}(v)]/\hbar$ becomes the Poisson bracket
$\{t_{ij}(u), t_{kl}(v)\}$.  The right side
$(t_{kj}(u)\,t_{il}(v) - t_{kj}(v)\,t_{il}(u))/(u-v)$
is already independent of $\hbar$ (the generators $t_{ij}(u)$
have classical limit $t_{ij}^{\mathrm{cl}}(u)$).
\end{proof}

\begin{proposition}[Sklyanin bracket satisfies Jacobi;
\ClaimStatusProvedHere]
\label{prop:sklyanin-jacobi}
The bracket~\eqref{eq:sklyanin-bracket} satisfies the Poisson
Jacobi identity.
\end{proposition}

\begin{proof}
The Jacobi identity for~\eqref{eq:sklyanin-bracket} is
equivalent to the classical Yang--Baxter equation (CYBE) for
$r(u) = P/u$:
$[r_{12}(u{-}v),\, r_{13}(u{-}w)]
+ [r_{12}(u{-}v),\, r_{23}(v{-}w)]
+ [r_{13}(u{-}w),\, r_{23}(v{-}w)] = 0$.
This is verified by direct computation: each commutator
produces terms of the form $P_{ij}P_{kl}/(u_a - u_b)(u_c - u_d)$,
and the six terms cancel pairwise by the permutation identity
$P_{12}P_{13} + P_{13}P_{23} + P_{23}P_{12}
= P_{12}P_{23} + P_{23}P_{13} + P_{13}P_{12}$.
Alternatively: this is the classical limit of the quantum
Yang--Baxter equation (Theorem~\ref{thm:YBE}), which is
proved in~\S\ref{subsec:YBE-proof}.
\end{proof}

\begin{computation}[Heisenberg verification;
\ClaimStatusProvedHere]
\label{comp:heisenberg-sklyanin}
For $\fg = \mathfrak{u}(1)$: $t(u) = J(u)$,
$P = 1$, and
$\{J(u), J(v)\} = (J(u) J(v) - J(v) J(u))/(u-v) = 0$
(the classical Heisenberg is abelian).  The Poisson bracket on
modes $J_r$ comes from the second-order pole of the
\emph{quantum} relation~\eqref{eq:heisenberg-spectral-OPE}:
$\{J_m, J_n\}_{\mathrm{qu}} = -k\,m\,\delta_{m+n,0}$.
This is the standard Heisenberg PVA at level $-k$, confirming
that $\cA^!_{\mathrm{line}} \simeq \widehat{\mathfrak{u}(1)}_{-k}$.
\end{computation}

\subsubsection*{The $A_\infty$ structure: what is proved and
what is open}

The $E_1$-chiral $A_\infty$ structure on $Y_\hbar(\fg)$ is
the central open question.  We state precisely what is known.

The Yangian $Y_\hbar(\fg)$ is a \emph{filtered} algebra:
the PBW filtration by total mode number $|t_{ij}^{(r)}| = r$
gives $\gr Y_\hbar(\fg) \simeq U(\fg[t])$.

\begin{proposition}[Associated graded is formal;
\ClaimStatusProvedHere]
\label{prop:gr-yangian-formal}
The associated graded $U(\fg[t])$ is Koszul as a graded
algebra.  Its bar complex is generated in degrees $1$ and $2$:
the generators are $\fg[t]$ in degree~$1$, and the relations
$xy - yx - [x,y] = 0$ in degree~$2$.  The $A_\infty$ structure
on $U(\fg[t])$ is formal: $m_k = 0$ for $k \ge 3$.
\end{proposition}

\begin{proof}
$U(\fg[t]) = U(\fg) \otimes \C[t]$ as a filtered algebra.
Both $U(\fg)$ (Priddy) and $\C[t]$ (trivially) are Koszul, and
the tensor product of Koszul algebras is Koszul
(Polishchuk--Positselski, Theorem~2.3.1).
\end{proof}

\begin{question}[$E_1$-chiral $A_\infty$ of $Y_\hbar(\fg)$;
\ClaimStatusOpen]
\label{q:yangian-Ainfty}
\index{Yangian!$A_\infty$ structure}
Is the full $Y_\hbar(\fg)$ (not just its associated graded)
formal as an $E_1$-chiral algebra?  The RTT
relation~\eqref{eq:FRT-components} is quadratic in the
generating functions $t_{ij}(u)$, but the mode-expanded
relations~\eqref{eq:yangian-relation-modes} mix all mode
numbers.  The algebra is a non-trivial filtered deformation
of a Koszul algebra.

Two scenarios:
\begin{enumerate}[label=\textup{(\alph*)}]
\item \emph{Filtered formality:}
  There exists a filtered $A_\infty$ quasi-isomorphism
  $Y_\hbar(\fg) \xrightarrow{\sim} (U(\fg[t]), m_2)$
  that is the identity on the associated graded.  This would
  mean $Y_\hbar(\fg)$ is ``formal up to filtration'': the
  higher $m_k$ vanish on $\gr$ but may be non-zero on the
  filtered level.
\item \emph{Genuine non-formality:}
  No such quasi-isomorphism exists, and the $A_\infty$
  operations $m_k$ for $k \ge 3$ are genuinely non-trivial
  (not killed by any change of generators).
\end{enumerate}
The answer depends on the vanishing of specific deformation
classes in the Harrison cohomology of $U(\fg[t])$ with
coefficients in the adjoint bimodule, twisted by the
$\hbar$-deformation.  This is not computed in either volume.
\end{question}

%% ----------------------------------------------------------------
%% Toward filtered formality: the Harrison cohomology computation
%% ----------------------------------------------------------------

\subsubsection*{Toward filtered formality: the Harrison cohomology computation}

We now give a first-principles attack on
Question~\ref{q:yangian-Ainfty}, providing strong evidence for
scenario~(a).  The key input is the Koszul property of
$U(\fg[t])$ and the K\"unneth decomposition for Harrison
cohomology.

\begin{proposition}[Hochschild decomposition for
{$U(\fg[t])$}; \ClaimStatusProvedHere]
\label{prop:HH-Ugt-decomposition}
\index{Hochschild cohomology!of $U(\fg[t])$}
The Hochschild cohomology of $U(\fg[t])$ with adjoint
coefficients decomposes as
\[
  \HH^n\bigl(U(\fg[t]),\, U(\fg[t])\bigr)
  \;\cong\;
  \bigoplus_{p+q=n}
  \Lambda^p(\fg[t]) \otimes \Sym^q(\fg[t]).
\]
\end{proposition}

\begin{proof}
Write $U(\fg[t]) = U(\fg) \otimes \bC[t]$.  Both tensor
factors are Koszul: $U(\fg)$ is Koszul by the
Poincar\'e--Birkhoff--Witt theorem (the Koszul resolution is
the Chevalley--Eilenberg complex), and $\bC[t]$ is the
symmetric algebra on a one-dimensional space, hence trivially
Koszul.  For a Koszul algebra $A$, the canonical Koszul
resolution provides a free $A^e$-resolution of $A$ whose
$\Hom$-complex computes Hochschild cohomology.  Applying this
to $U(\fg[t])$ and using the Hochschild--Kostant--Rosenberg
theorem for the smooth algebra $U(\fg[t])$ (which, as an
enveloping algebra of a free $\bC[t]$-Lie algebra, has
homological regularity), one obtains the stated
decomposition.  The exterior powers $\Lambda^p(\fg[t])$
arise from the antisymmetric Hochschild cochains (the
polyvector fields), while $\Sym^q(\fg[t])$ arises from the
polynomial functions on the formal neighborhood of the
diagonal.  The Harrison cohomology
$\Har^n\bigl(U(\fg[t]),\, U(\fg[t])\bigr)$ sits inside as
the indecomposable summand: the subspace of cochains that
vanish on all non-trivial shuffles.
\end{proof}

\begin{proposition}[Deformation class and obstruction locus;
\ClaimStatusProvedHere]
\label{prop:yangian-deformation-obstruction}
\index{Yangian!deformation class}
\index{Harrison cohomology!obstruction to formality}
The Yangian deformation class
$[Y_\hbar(\fg)] \in \HH^2\bigl(U(\fg[t]),\, U(\fg[t])\bigr)$
is the class of the RTT deformation: in the decomposition of
Proposition~\textup{\ref{prop:HH-Ugt-decomposition}}, it
lies in the $\Lambda^2(\fg[t])$ summand and is represented by
the classical $r$-matrix $r(u-v) \in \fg \otimes \fg \otimes
\bC(u-v)$, expanded in the mode basis.

The obstruction to extending this deformation to a filtered
$A_\infty$ quasi-isomorphism
$Y_\hbar(\fg) \xrightarrow{\sim} (U(\fg[t]), m_2)$ lies in
the Harrison cohomology group
$\Har^3\bigl(U(\fg[t]),\, U(\fg[t])\bigr)$.
\end{proposition}

\begin{proof}
The first statement follows from the standard identification
of infinitesimal deformations of an associative algebra $A$
with $\HH^2(A,A)$: the Yangian multiplication
$m_\hbar = m_0 + \hbar\, \mu_1 + O(\hbar^2)$ defines a
Hochschild $2$-cocycle $\mu_1$, and the RTT
presentation~\eqref{eq:FRT-components} identifies $\mu_1$
with the $r$-matrix class.

For the second statement: a filtered $A_\infty$
quasi-isomorphism $f = (f_1, f_2, \ldots)$ with $f_1 = \id$
is constructed inductively in powers of $\hbar$.  At each
order, the obstruction to extending $f$ to the next order is
a Harrison $3$-cocycle (not merely a Hochschild $3$-cocycle),
since the $A_\infty$ relations constrain the cochains to be
indecomposable; the shuffle-trivial part is absorbed by
lower-order corrections to $f_k$ for $k \ge 2$.
The obstruction therefore lives in
$\Har^3\bigl(U(\fg[t]),\, U(\fg[t])\bigr)$.
\end{proof}

\begin{lemma}[Vanishing of the leading Harrison obstruction;
\ClaimStatusProvedHere]
\label{lem:harrison-leading-vanishing}
\index{Harrison cohomology!K\"unneth decomposition}
By the tensor decomposition
$U(\fg[t]) = U(\fg) \otimes \bC[t]$ and the K\"unneth
formula for Harrison cohomology of tensor products of
commutative algebras, one has
\begin{equation}
\label{eq:harrison-kunneth}
  \Har^3\bigl(U(\fg[t])\bigr)
  \;\cong\;
  \bigoplus_{i+j=3}
  \Har^i\bigl(U(\fg)\bigr) \otimes
  \Har^j\bigl(\bC[t]\bigr).
\end{equation}
The leading term $\Har^3(U(\fg)) \otimes \Har^0(\bC[t])$
vanishes: since $U(\fg)$ is Koszul, its Harrison obstruction
group is trivial by the theorem of Polishchuk--Positselski
\textup{\cite{polishchuk2005quadratic}} that Koszul algebras
have vanishing higher Harrison cohomology
$\Har^n = 0$ for $n \ge 3$.
\end{lemma}

\begin{proof}
The K\"unneth formula for Harrison cohomology of a tensor
product $A \otimes B$ of augmented commutative algebras is
the standard shuffle K\"unneth
(cf.~\cite{loday1992cyclic}): Harrison cochains on
$A \otimes B$ decompose along the tensor bigrading, with no
higher Tor correction since we work over a field.  The
summands in~\eqref{eq:harrison-kunneth} are:
\begin{align*}
  (i,j) = (3,0)&\colon\quad
    \Har^3(U(\fg)) \otimes \bC = 0,\\
  (i,j) = (2,1)&\colon\quad
    \Har^2(U(\fg)) \otimes \Har^1(\bC[t]),\\
  (i,j) = (1,2)&\colon\quad
    \Har^1(U(\fg)) \otimes \Har^2(\bC[t]),\\
  (i,j) = (0,3)&\colon\quad
    \bC \otimes \Har^3(\bC[t]) = 0,
\end{align*}
where $\Har^3(\bC[t]) = 0$ since $\bC[t]$ is a polynomial
algebra (hence smooth, hence Harrison-formal).  The
vanishing $\Har^3(U(\fg)) = 0$ is the content of the
Polishchuk--Positselski result: for a Koszul algebra, the
Harrison complex is quasi-isomorphic to the space of Koszul
generators, which is concentrated in degrees $\le 2$.
\end{proof}

\begin{remark}[Evidence for filtered formality]
\label{rem:filtered-formality-evidence}
Lemma~\ref{lem:harrison-leading-vanishing} kills the leading
obstruction to filtered formality, which is the term
requiring no polynomial dependence on the spectral parameter
$t$.  The surviving terms $\Har^2(U(\fg)) \otimes
\Har^1(\bC[t])$ and $\Har^1(U(\fg)) \otimes
\Har^2(\bC[t])$ involve \emph{mixed} Harrison classes with
non-trivial $t$-dependence.  These require a separate
vanishing argument; one expects them to be killed by the
$\fg$-equivariance constraint on the deformation (the
Yangian is a $\fg$-equivariant deformation of $U(\fg[t])$,
and the mixed terms are not $\fg$-invariant at the relevant
weights).

There is also an independent line of evidence: the
Etingof--Kazhdan quantization theorem
\cite{etingof1996quantization} establishes formality of the
quantum group $U_q(\fg)$ as a deformation of $U(\fg)$.  The
Yangian $Y_\hbar(\fg)$ is the rational degeneration of the
quantum affine algebra $U_q(\widehat{\fg})$, and the
Etingof--Kazhdan formality machinery (itself
cohomological in nature, controlling obstructions via Drinfeld
associators) should survive the rational limit
$q = e^{\hbar/u} \to 1 + \hbar/u$.  A complete proof would
require verifying that the relevant pro-nilpotent completion
is compatible with the rational degeneration, which we leave
to future work.
\end{remark}

\begin{remark}[The Steinberg perspective on filtered formality]
\label{rem:steinberg-filtered-formality}
\index{Steinberg variety!filtered formality}
Recall that $Y_\hbar(\fg)$ acts on $\cO$-modules on the
Steinberg variety $Z = \widetilde{\cN} \times_{\cN}
\widetilde{\cN}$, just as $U(\fg[t])$ acts on the undeformed
correspondence in the factorisation-quantum-group picture above.
Filtered formality of $Y_\hbar(\fg) \xrightarrow{\sim}
(U(\fg[t]), m_2)$ is equivalent to a formal equivalence
between the deformed and undeformed Steinberg categories
$D^b\mathrm{Coh}(Z_\hbar) \simeq D^b\mathrm{Coh}(Z)$ as
filtered dg-categories.  Under the localization dictionary,
the Harrison obstruction in
$\Har^3(U(\fg[t]))$ maps to the obstruction to lifting the
identity functor on $D^b\mathrm{Coh}(Z)$ to a filtered
equivalence with $D^b\mathrm{Coh}(Z_\hbar)$.  The vanishing
of the leading Harrison term
(Lemma~\ref{lem:harrison-leading-vanishing}) thus translates
geometrically: the undeformed Steinberg variety has no
intrinsic third-order obstruction to formal deformation, and
the remaining obstructions are controlled by spectral-parameter-dependent classes that mix the
geometric and loop directions.
\end{remark}

\begin{remark}[Virasoro: non-formality on both sides]
\label{rem:yangian-virasoro-nonformality}
For the Virasoro algebra $\mathrm{Vir}_c$ at generic $c$,
the Koszul dual $\mathrm{Vir}_{26-c}$ is \emph{not} formal:
$m_k \ne 0$ for all $k \ge 3$
(Section~\ref{subsec:gravity-shadow-tower}).  The Virasoro
has no Yangian simplification: both sides of the Koszul
duality are genuinely non-formal, and the infinite-depth
shadow tower persists on the line-operator side.
\end{remark}

\subsubsection*{The bar-cobar involution}

The Koszul involution $(-)^{!!} \simeq \id$ at the
$E_1$-chiral level gives
$\Omega^{E_1\text{-ch}}(B^{E_1\text{-ch}}(\cA)) \simeq \cA$.
For $\cA = \widehat{\fg}_k$: the ordered bar produces
$Y_\hbar(\fg)$, and the cobar recovers $\widehat{\fg}_k$.

The structures exchange:
\begin{center}
\small
\renewcommand{\arraystretch}{1.15}
\begin{tabular}{lll}
\textbf{Boundary $\cA$} &
  \textbf{Koszul dual $\cA^!$} &
  \textbf{Mechanism} \\
\hline
Current OPE $J^a(z)J^b(w)$ &
  FRT relation~\eqref{eq:FRT-components} &
  bar $\leftrightarrow$ cobar \\[2pt]
Poles in $z - w$ &
  Poles in $u - v$ &
  holomorphic $\leftrightarrow$ spectral \\[2pt]
$\lambda$-bracket &
  Sklyanin bracket~\eqref{eq:sklyanin-bracket} &
  PVA $\leftrightarrow$ PVA \\[2pt]
$r(z) = \Omega/z$ &
  $r_{\mathrm{cl}}(u) = P/u$ &
  same Casimir \\[2pt]
$\kappa = \dim(\fg)(k{+}h^\vee)/(2h^\vee)$ &
  $\kappa^! = -\kappa$ &
  complementarity \\[2pt]
$\SCchtop$-algebra &
  $\SCchtop$-algebra &
  two-color involution (Thm~\ref{thm:two-color-master})
\end{tabular}
\end{center}

\begin{computation}[Heisenberg bar-cobar;
\ClaimStatusProvedHere]
\label{comp:heisenberg-bar-cobar}
For $\cA = \widehat{\mathfrak{u}(1)}_k$: the boundary has
$J(z)J(w) \sim k/(z-w)^2$, and the open-colour dual
$\cA^!_{\mathrm{line}} \simeq \widehat{\mathfrak{u}(1)}_{-k}$ has
$J(u)J(v) \sim -k/(u-v)^2$.  The ordered bar complex of
$\cA^!_{\mathrm{line}}$ should recover $\cA$.  Indeed: the
ordered bar differential on
$[s^{-1}J \,|\, s^{-1}J]$ extracts the
residue of $-k/(u-v)^2$, which is $k\,\partial$, the level-$k$
Heisenberg OPE.  The level flips sign twice under the double
bar and returns to $k$, recovering the boundary algebra~$\cA$.
\end{computation}

\subsubsection*{Genus tower from complementarity}

Koszul complementarity (Vol~I, Theorem~C, the Lagrangian splitting of genus-$g$ cohomology) gives
\begin{equation}\label{eq:yangian-kappa}
\kappa(\cA^!)
\;=\;
-\kappa(\cA).
\end{equation}
For $\cA = \widehat{\fg}_k$:
$\kappa = \dim(\fg)(k{+}h^\vee)/(2h^\vee)$, so
$\kappa(\cA^!) = -\kappa(\cA)$.  The complementarity
potential $F_1(\cA) + F_1(\cA^!) = 0$ vanishes: the bulk
theory carries no net genus-$1$ anomaly.

The all-genus generating function:
\[
\sum_{g \ge 1} F_g(Y_\hbar(\fg))\,x^{2g}
\;=\;
-\kappa(\widehat{\fg}_k)\bigl(\hat{A}(ix) - 1\bigr).
\]

\begin{computation}[Heisenberg genus tower;
\ClaimStatusProvedHere]
\label{comp:heisenberg-genus}
For $\cA = \widehat{\mathfrak{u}(1)}_k$:
$\kappa = k$, $\kappa^! = -k$.  The genus-$1$ free energies
are $F_1(\cA) = k/24$ and $F_1(\cA^!) = -k/24$.  Their sum
vanishes: the bulk $\mathfrak{u}(1)$ theory has no
genus-$1$ anomaly.  This matches the direct computation:
the bulk partition function on $T^2$ is
$Z_{\mathrm{bulk}} = |\eta(\tau)|^{-2}$, which is modular
invariant for all $k$.
\end{computation}

\begin{remark}[What remains open]
\label{rem:yangian-E1-chiral-open}
Four problems are open:
\begin{enumerate}[label=\textup{(\roman*)}]
\item The $A_\infty$ formality question
  (Question~\ref{q:yangian-Ainfty}).
\item Construction of a ``Yangian vertex algebra'': a vertex
  algebra-like object on $Y_\hbar(\fg)$ encoding the spectral
  OPE, state-field correspondence, and translation.  The
  obstruction is that the Yangian is filtered by mode number, not
  graded by conformal weight.  Theorem~\ref{thm:dual-sc-algebra}
  resolves the structural question (the open-colour dual
  $\cA^!_{\mathrm{line}}$ carries a canonical
  $\SCchtop$-algebra structure encoding both the Sklyanin
  $\lambda$-bracket (closed colour) and the Yangian product
  (open colour)), but the explicit ``vertex algebra''
  presentation remains a presentation problem.
\item Quantization of the Sklyanin PVA
  \eqref{eq:sklyanin-bracket} via the modular programme
  (Chapter~\ref{ch:modular-pva-quantization}).
\item The $E_1$-chiral $A_\infty$ structure on
  non-formal Koszul duals ($\mathrm{Vir}_{26-c}$,
  $\mathcal{W}$-algebras at dual parameters).
\end{enumerate}
\end{remark}


%%%%%%%%%%%%%%%%%%%%%%%%%%%%%%%%%%%%%%%%%%%%%%%%%%%%%%%%%%%%%%%%%%%%%%%%%%%%%
% TWO-COLOR KOSZUL DUALITY
%%%%%%%%%%%%%%%%%%%%%%%%%%%%%%%%%%%%%%%%%%%%%%%%%%%%%%%%%%%%%%%%%%%%%%%%%%%%%

\subsection{Two-color Koszul duality}
\label{subsec:two-color-koszul-duality}
\index{two-color Koszul duality|textbf}
\index{Swiss-cheese operad!two-color Koszul duality}
\index{Koszul duality!two-color}

The bar complex $\barB(\cA)$ carries two compatible structures:
the bar differential $d_{\barB}$ (from OPE residues on
$\FM_k(\C)$) and the deconcatenation coproduct $\Delta$
(from ordered configurations on $\Conf_k^{<}(\R)$).
Together these make $\barB(\cA)$ an $\SCchtop$-coalgebra.
The homotopy-Koszulity of $\SCchtop$
(Theorem~\ref{thm:homotopy-Koszul}) and the associated graded
splitting (Proposition~\ref{prop:gr-chiral}) produce a
\emph{simultaneous} Koszul duality in both colors.  This
subsection extracts the master theorem.

\subsubsection*{The two-color Koszul datum}

\begin{definition}[Two-color Koszul datum]
\label{def:two-color-koszul-datum}
\index{two-color Koszul datum}
Let $\cA$ be a logarithmic $\SCchtop$-algebra
\textup{(}Definition~\textup{\ref{def:log-SC-algebra})} and
let $\cA^!_{\mathrm{ch}}$ and $\cA^!_{\mathrm{line}}$ denote
its closed- and open-colour Koszul duals.  The \emph{two-color
Koszul datum} is the septuple
\[
\bigl(\cA,\; \cA^!_{\mathrm{ch}},\; \cA^!_{\mathrm{line}},\;
\barBch(\cA),\; \barB^{\mathrm{ord}}(\cA),\;
\tau_{\mathrm{ch}},\; \tau_{\mathrm{ord}}\bigr),
\]
where:
\begin{enumerate}[label=\textup{(\roman*)}]
\item $\barBch(\cA)$ is the \emph{chiral bar complex}
  (Definition~\ref{def:chiral_bar}): the bar construction with
  differential from $\FM_k(\C)$-integrals.  This is the
  closed-color sector of the $\SCchtop$-bar complex.
\item $\barB^{\mathrm{ord}}(\cA)$ is the \emph{ordered bar
  complex}: the cofree coalgebra $T^c(s^{-1}\bar{\cA})$ with
  deconcatenation coproduct from $\Conf_k^{<}(\R)$.  This is
  the open-color sector.
\item $\tau_{\mathrm{ch}} \colon \barBch(\cA) \to
  \cA^!_{\mathrm{ch}}$ and
  $\tau_{\mathrm{ord}} \colon \barB^{\mathrm{ord}}(\cA) \to
  \cA^!_{\mathrm{line}}$ are the closed-color and open-color
  universal twisting morphisms, respectively.
\end{enumerate}
The two bar complexes are \emph{not independent}: they are the
two projections of the single $\SCchtop$-bar complex
$\barB^{\SCchtop}(\cA)$.
\end{definition}

\subsubsection*{The master theorem}

\begin{theorem}[Two-color Koszul duality;
\ClaimStatusProvedHere]
\label{thm:two-color-master}
\index{two-color Koszul duality!master theorem|textbf}
Let $\cA$ be a chirally Koszul logarithmic $\SCchtop$-algebra
\textup{(}Definition~\textup{\ref{def:log-SC-algebra})}.
The homotopy-Koszulity of $\SCchtop$
simultaneously produces:
\begin{enumerate}[label=\textup{(\roman*)}]
\item \emph{Closed-color duality.}
  A Quillen equivalence between chiral $\cA$-modules and
  chiral $\cA^!_{\mathrm{ch}}$-comodules, recovering
  Theorems~A and~B of~Vol~I\@.  The bar functor
  $\barBch \colon \cA\text{-mod}_{\mathrm{ch}} \rightleftarrows
  \cA^!_{\mathrm{ch}}\text{-comod}_{\mathrm{ch}} \cocolon \Omegach$ is
  the closed-color bar-cobar adjunction.
\item \emph{Open-color duality.}
  A Quillen equivalence between $E_1$ $\cA$-modules and $E_1$
  $\cA^!_{\mathrm{line}}$-comodules, identifying
  $\mathcal{C}_{\mathrm{line}} \simeq \cA^!_{\mathrm{line}}\text{-mod}$
  (Theorem~\ref{thm:lines_as_modules}).  The bar functor
  $\barB^{\mathrm{ord}} \colon \cA\text{-mod}_{E_1}
  \rightleftarrows
  \cA^!_{\mathrm{line}}\text{-comod}_{E_1} \cocolon
  \Omega^{\mathrm{ord}}$
  is the open-color bar-cobar adjunction.
\item \emph{Interaction.}
  The holomorphic weight filtration on $\SCchtop$ induces a
  spectral sequence $E_r$ converging to
  $H^\ast(\barB^{\SCchtop}(\cA))$ with:
  \begin{itemize}
  \item $E_0$-page: the two colors decoupled
    (Proposition~\ref{prop:gr-chiral});
  \item $d_1$-differentials encoding the classical $r$-matrix
    $r_0/z$ (leading residue of the $R$-matrix);
  \item $d_r$-differentials encoding the $r$-th OPE pole
    coefficient of $R(z)$;
  \item convergence to the full $\SCchtop$-coalgebra structure
    on~$\barB(\cA)$.
  \end{itemize}
\item \emph{Involution.}
  $(-)^{!!} \simeq \id$ at the $\SCchtop$ level:
  $\Omega^{\SCchtop}(\barB^{\SCchtop}(\cA)) \simeq \cA$
  as $\SCchtop$-algebras.  This upgrades the separate
  closed-color and open-color involutions to a single
  two-color involution.
\end{enumerate}
\end{theorem}

\begin{proof}
By Theorem~\ref{thm:homotopy-Koszul}, the operad $\SCchtop$
is homotopy-Koszul.  The general theory of homotopy-Koszul
operads (Loday--Vallette \cite{LV12}, Chapter~10; generalized
to two-colored operads by Vallette~\cite{Val07}) gives a Quillen
equivalence between $\SCchtop$-algebras and
$(\SCchtop)^!$-coalgebras via the operadic bar-cobar
adjunction.  We now project to each color.

\medskip
\noindent\textbf{(i) Closed color.}
The closed-color projection of $\barB^{\SCchtop}(\cA)$ retains
only the $\FM_k(\C)$-labeled operations, giving $\barBch(\cA)$.
The Quillen equivalence restricted to closed-color modules is
the chiral bar-cobar adjunction, which is Theorems~A and~B
of Vol~I specialized to genus~$0$ and the standard curve
$X = \C$.  The key input: the closed-color factor of
$\mathrm{gr}\;\SCchtop$ is the chiral operad
$\mathsf{P}_{\mathrm{ch}}$, which is Koszul
(Francis--Gaitsgory~\cite{FG12}); this is the projection
landing on the closed-color dual~$\cA^!_{\mathrm{ch}}$.

\medskip
\noindent\textbf{(ii) Open color.}
The open-color projection retains only the
$E_1(m)$-labeled operations, giving the ordered bar complex
$\barB^{\mathrm{ord}}(\cA)$.  The Quillen equivalence
restricted to open-color modules identifies
$\cA^!_{\mathrm{line}}$-modules with line operators, recovering
Theorem~\ref{thm:lines_as_modules}.

\medskip
\noindent\textbf{(iii) Interaction.}
The holomorphic weight filtration $F^\bullet$ on $\SCchtop$ is
the filtration by pole order on $\FM_k(\C)$.  By
Proposition~\ref{prop:gr-chiral},
$\mathrm{gr}^F\,\SCchtop \cong \mathsf{P}_{\mathrm{ch}}
\amalg E_1$ with no mixed operations.  The spectral sequence of
the filtered complex $(\barB^{\SCchtop}(\cA), F^\bullet)$
therefore has $E_0 = \barBch(\cA) \otimes
\barB^{\mathrm{ord}}(\cA)$ (decoupled).

The $d_1$ differential is the leading mixed-color correction:
it acts on bar degree~$2$ elements $[a_1 | a_2]$ by extracting
the simple pole $\Res_{z=0} \{a_{1\,\lambda}\, a_2\}/z$,
which is exactly the classical $r$-matrix $r_0 = \Omega/z$
(Proposition~\ref{prop:field-theory-r}).  At bar degree~$k$,
$d_r$ encodes the order-$r$ pole in the OPE:
$\Res_{z=0} z^{r-1} R(z)$ restricted to $k$-fold bar elements.

Convergence: the filtration is bounded below
(pole order~$\ge 0$) and exhaustive ($\FM_k(\C)$ has finite
pole orders for finite $k$), so the spectral sequence converges
by the standard bounded-filtration theorem.

\medskip
\noindent\textbf{(iv) Involution.}
The two-color bar-cobar composite
$\Omega^{\SCchtop} \circ \barB^{\SCchtop}$ is the identity on
the homotopy category of $\SCchtop$-algebras by the general
Koszul duality involution (Loday--Vallette, Theorem~10.1.13).
Since the closed-color and open-color projections of this
composite recover the separate involutions $\Omegach \circ
\barBch \simeq \id$ and
$\Omega^{\mathrm{ord}} \circ \barB^{\mathrm{ord}} \simeq
\id$, the two-color involution refines both.
\end{proof}

\subsubsection*{The spectral parameter from the closed color}

\begin{theorem}[Spectral parameter from the closed color;
\ClaimStatusProvedHere]
\label{thm:spectral-parameter-from-closed-color}
\index{spectral parameter!closed-color origin}
\index{R-matrix!two-color origin}
The spectral parameter $z$ in $r(z)$ is the coordinate on
$\FM_2(\C)$ in the closed-color factor of $\SCchtop$.
\begin{enumerate}[label=\textup{(\roman*)}]
\item The $E_1$ structure sees ordering on $\R_t$: the
  deconcatenation coproduct $\Delta$ on
  $\barB^{\mathrm{ord}}(\cA)$ is a map of graded vector spaces,
  independent of any spectral parameter.
\item The chiral structure sees holomorphic separation
  $z = z_1 - z_2$ on $\FM_2(\C)$: the bar differential
  $d_{\barB}$ on $\barBch(\cA)$ depends meromorphically on $z$.
\item The $R$-matrix lives in the \emph{mixed sector}:
  holomorphic weight $\ge 1$ (closed color contributes poles
  in $z$) and open-color degree~$2$ (acts on pairs of line
  operators via $\Delta$).
\end{enumerate}
\end{theorem}

\begin{proof}
Item~(i): the deconcatenation coproduct
$\Delta \colon T^c(V) \to T^c(V) \otimes T^c(V)$ is the
standard unshuffle coproduct on the cofree coalgebra, which is
defined purely combinatorially (Loday--Vallette, \S1.2).  No
geometric data enters; $\Delta$ is the identity of the
open-color $E_1$-coalgebra.

Item~(ii): the bar differential $d_{\barB}$ is constructed
from configuration integrals over $\FM_k(\C)$
(Sections~\ref{sec:BV_construction}--\ref{sec:FM_calculus}).
The fiber of $\FM_2(\C)$ over the diagonal is parametrized by
the tangent direction, whose coordinate is $z = z_1 - z_2$.
The OPE $a_1(z_1) a_2(z_2) \sim \sum_n C_n (z_1-z_2)^{-n-1}$
makes $d_{\barB}$ meromorphic in $z$.

Item~(iii): in the holomorphic weight filtration, an element of
$F^p$ has pole order $\ge p$ on $\FM_k(\C)$.  The identity
(pole order~$0$) is open-color-only; the $R$-matrix $r(z)$
has pole order~$1$ (the simple pole $r_0/z$) and acts on pairs
(bar degree~$2$).  Hence $r(z) \in F^1 \cap
\barB^2(\cA) \otimes \barB^2(\cA)$: mixed-sector data.
\end{proof}

\subsubsection*{The spectral sequence for the $R$-matrix}

\begin{proposition}[Spectral sequence for the $R$-matrix;
\ClaimStatusProvedHere]
\label{prop:spectral-sequence-r-matrix}
\index{spectral sequence!R-matrix}
\index{R-matrix!spectral sequence}
The holomorphic weight filtration on $\barB^{\SCchtop}(\cA)$
induces a spectral sequence $\{E_r, d_r\}_{r \ge 0}$:
\begin{enumerate}[label=\textup{(\roman*)}]
\item $E_0 = \barBch(\cA) \otimes \barB^{\mathrm{ord}}(\cA)$
  (the two colors decoupled, by
  Proposition~\ref{prop:gr-chiral}).
\item $d_1 \colon E_0^{p,q} \to E_0^{p+1,q}$ encodes the
  leading $R$-matrix coefficient: on bar degree~$2$ elements,
  $d_1([a_1 | a_2]) = r_0(a_1, a_2)/z$ where
  $r_0 = \sum_{ij} \Omega^{ij}\, e_i \otimes e_j$ is the
  Casimir (Proposition~\ref{prop:field-theory-r}).
\item At the $r$-th page, $d_r$ encodes the coefficient of
  $(z_1-z_2)^{-r-1}$ in the OPE: the $r$-th order correction
  to the $R$-matrix.
\item The spectral sequence converges:
  $E_\infty \cong H^\ast(\barB^{\SCchtop}(\cA))$.
\end{enumerate}
\end{proposition}

\begin{proof}
The spectral sequence is that of the filtered chain complex
$(\barB^{\SCchtop}(\cA), F^\bullet, d^{\SCchtop})$ where
$F^p$ consists of elements with holomorphic pole order $\ge p$.
The differential $d^{\SCchtop}$ splits as
$d^{\SCchtop} = d_0 + d_1 + d_2 + \cdots$ where $d_r$ raises
pole order by exactly $r$.  At $E_0$: $d_0$ is the sum of the
closed-color and open-color differentials, which are
independent by Proposition~\ref{prop:gr-chiral}.  The
$d_1$-differential is the leading mixed-color contribution:
from the simple pole in the OPE $a_1(z) a_2(0) \sim
C_0/(z) + \mathrm{regular}$, this is $r_0/z$.

Convergence: the filtration is bounded below ($F^0 =
\barB^{\SCchtop}(\cA)$, $F^p = 0$ for $p$ exceeding the
maximal pole order on $\FM_k(\C)$ at fixed bar degree~$k$)
and complete.  The spectral sequence of a bounded-below
exhaustive filtration on a cochain complex converges
(Weibel~\cite{Wei94}, Theorem~5.5.1).
\end{proof}

\begin{corollary}[$\Ainf$ Yang--Baxter as convergence;
\ClaimStatusProvedHere]
\label{cor:ainfty-ybe-as-convergence}
\index{Yang--Baxter equation!$A_\infty$}
\index{Yang--Baxter equation!spectral sequence convergence}
The $\Ainf$ Yang--Baxter equation
(Theorem~\ref{thm:spectral_R_YBE}) is the convergence condition
of the spectral sequence $E_r$ at bar degree~$3$: the tower of
differentials $d_1, d_2, \ldots$ acting on three-fold bar
elements $[a_1 | a_2 | a_3]$ converges if and only if the
full $R$-matrix $R(z)$ satisfies the quantum Yang--Baxter
equation on triple tensor products.
\end{corollary}

\begin{proof}
At bar degree~$3$, the $E_r$-page carries the data of three
pairwise interactions $(12)$, $(13)$, $(23)$ at holomorphic
separations $z_{12}$, $z_{13}$, $z_{23}$.  The $d_r$
differentials at this bar degree are:
\begin{align*}
d_1([a_1 | a_2 | a_3])
&= r_0^{(12)}/z_{12} + r_0^{(13)}/z_{13}
   + r_0^{(23)}/z_{23}, \\
d_r([a_1 | a_2 | a_3])
&= \text{order-$r$ pole corrections to each pair}.
\end{align*}
The condition $(d^{\SCchtop})^2 = 0$ on bar degree~$3$
elements reads, at each pole order:
\begin{itemize}
\item Order $2$: $[r_0^{(12)}, r_0^{(13)}]
  + [r_0^{(12)}, r_0^{(23)}]
  + [r_0^{(13)}, r_0^{(23)}] = 0$
  \quad (the classical YBE);
\item Order $r+1$: the correction from $d_r$ compensates
  the failure of $d_1^2 + \cdots + d_{r-1}^2$ at the previous
  stage.
\end{itemize}
The tower converges if and only if these corrections
sum to zero at all orders, which is exactly the quantum YBE
for $R(z) = \id + \hh\,r_0/z + O(z^{-2})$.
\end{proof}

\subsubsection*{The $\SCchtop$-algebra structure on $\cA^!_{\mathrm{line}}$}

\begin{theorem}[Dual $\SCchtop$-algebra;
\ClaimStatusProvedHere]
\label{thm:dual-sc-algebra}
\index{Swiss-cheese operad!dual algebra structure}
\index{Koszul dual!$\SCchtop$-algebra}
\index{Sklyanin bracket!as closed-color cobar}
\index{Yangian!as open-color cobar}
The open-colour Koszul dual $\cA^!_{\mathrm{line}}$ inherits a canonical
$\SCchtop$-algebra structure:
\begin{enumerate}[label=\textup{(\roman*)}]
\item \emph{Closed color:} the coalgebra structure on
  $\barBch(\cA)$ induces, via Verdier duality on $\Ran(X)$,
  the Sklyanin $\lambda$-bracket on $\cA^!_{\mathrm{line}}$
  (equation~\eqref{eq:sklyanin-bracket}).
\item \emph{Open color:} the ordered cobar construction
  $\Omega^{\mathrm{ord}}(\barB^{\mathrm{ord}}(\cA))$ endows
  $\cA^!_{\mathrm{line}}$ with its Yangian product.
\item \emph{Compatibility:} $d_{\barB}$ is a coderivation
  of~$\Delta$ in $\barB^{\SCchtop}(\cA)$.  Dually, the
  Sklyanin bracket and the Yangian product satisfy the
  Swiss-cheese composition law: the $\lambda$-bracket is a
  derivation of the Yangian product up to homotopy controlled
  by the mixed-color operations of $\SCchtop$.
\end{enumerate}
The $\SCchtop$-algebra on $\cA^!_{\mathrm{line}}$ is the live
strict model of the two-color Koszul dual package relevant for
line operators and is Koszul dual to the
$\SCchtop$-coalgebra on $\barB^{\SCchtop}(\cA)$.
\end{theorem}

\begin{proof}
The operadic bar-cobar duality for homotopy-Koszul operads
gives $(\SCchtop)^! \simeq (\SCchtop)^{\vee}$
(Ginzburg--Kapranov~\cite{GK94}, generalized by Vallette
\cite{Val07} to the colored setting).  An
$\SCchtop$-coalgebra structure on $\barB(\cA)$ dualizes to an
$\SCchtop$-algebra structure on the strict model
$\cA^!_{\mathrm{line}}$ used on the live line/Yangian surface.
The two projections:

(i) The closed-color projection of the $\SCchtop$-algebra
structure maps are
$C_\ast(\FM_k(\C)) \otimes
(\cA^!_{\mathrm{line}})^{\otimes k}
\to \cA^!_{\mathrm{line}}$.
At $k = 2$, the $\FM_2(\C)$-integral with meromorphic kernel
gives the $\lambda$-bracket; the Sklyanin form
\eqref{eq:sklyanin-bracket} is the explicit residue computation
(Proposition~\ref{prop:field-theory-r}).

(ii) The open-color projection is
$E_1(m) \otimes (\cA^!_{\mathrm{line}})^{\otimes m}
\to \cA^!_{\mathrm{line}}$: the
$E_1$-algebra structure is the Yangian product, i.e.\ the
concatenation algebra structure dual to deconcatenation.

(iii) The mixed-color composition maps of $\SCchtop$ enforce
compatibility: the bar differential $d_{\barB}$, being a
coderivation of $\Delta$, dualizes to the statement that the
$\lambda$-bracket (from $d_{\barB}$) is a derivation of the
product (from $\Delta$) up to higher $\SCchtop$-homotopies.
\end{proof}

\begin{remark}[This is not an accident]
\label{rem:not-accidental}
\index{Yangian!not accidental}
The Yangian $Y_\hh(\fg)$ is \emph{not} ``an associative algebra
that also happens to carry a Poisson vertex algebra structure.''
It is an $\SCchtop$-algebra: the RTT relations encode the
closed color ($\FM_2(\C)$-integrals), the Yangian product
encodes the open color ($E_1$-composition), and the
$\lambda$-brackets are the closed-color OPE\@.  The Drinfeld
presentation separates these into ``Yangian generators'' (open)
and ``spectral parameter dependence'' (closed); the
$\SCchtop$-algebra structure reunifies them.
\end{remark}

\subsubsection*{Genus tower: asymmetry between the two colors}

\begin{theorem}[Genus tower asymmetry;
\ClaimStatusProvedHere]
\label{thm:genus-tower-asymmetry}
\index{genus tower!two-color asymmetry}
\index{curvature!closed-color only}
At genus $g \ge 1$, the two colors behave asymmetrically:
\begin{enumerate}[label=\textup{(\roman*)}]
\item \emph{Closed color is curved.}
  $\dfib^{\,2} = \kappa(\cA) \cdot \omega_g$
  (Theorem~\ref{thm:mc5-genus-one-bridge} at $g=1$;
  Vol~I, Theorem~C at all genera).  The bar complex
  $\barBch(\cA)$ on $\Sigma_g$ has non-vanishing curvature
  proportional to the modular characteristic $\kappa(\cA)$
  and the Arakelov form $\omega_g$.
\item \emph{Open color is flat.}
  The deconcatenation coproduct $\Delta$ is coassociative at
  \emph{all} genera: $(\Delta \otimes \id) \circ \Delta
  = (\id \otimes \Delta) \circ \Delta$ with no genus
  corrections.  The ordered bar complex
  $\barB^{\mathrm{ord}}(\cA)$ is a strict
  (uncurved) coalgebra at every genus.
\item \emph{Operadic origin.}
  The asymmetry is the ``no open-to-closed'' constraint of
  $\SCchtop$ (Definition~\ref{def:SC-colors}): closed-color
  operations depend on the complex structure of $\Sigma_g$
  via $\FM_k(\Sigma_g)$, which develops quasi-periodicity
  defects at $g \ge 1$.  Open-color operations
  depend only on $E_1(m)$, which is genus-independent since
  $\R$ is contractible at all genera.  The Arnold defect
  (quasi-periodic monodromy of $\FM_k(\Sigma_g)$) contributes
  only to the closed color.
\item \emph{Complementarity.}
  $\kappa(\cA) + \kappa(\cA^!) = 0$ for Kac--Moody and
  free-field algebras; more generally
  $\kappa + \kappa^! = \rho_\cA$ is a constant depending
  only on the algebra type
  (Vol~I, Theorem~C).  For the affine and free-field
  lineages ($\rho_\cA = 0$), the total
  $\SCchtop$-coalgebra is ``uncurved on average'': the
  closed-color curvatures of $\cA$ and $\cA^!$ cancel.
\end{enumerate}
\end{theorem}

\begin{proof}
(i) is Theorem~\ref{thm:mc5-genus-one-bridge} at genus~$1$,
extended to all genera by the genus tower of Vol~I\@.  The bar
differential on $\Sigma_g$ acquires curvature from the period
matrix: $d_{\barB}^2 = \kappa(\cA) \cdot \omega_g$ where
$\omega_g$ is the canonical Arakelov $(1,1)$-form
(equation~\eqref{eq:curved-R-fact}).

(ii): The deconcatenation coproduct $\Delta$ on the cofree
coalgebra $T^c(V)$ is defined purely combinatorially by
$(v_1 | \cdots | v_n) \mapsto
\sum_{i=0}^n (v_1 | \cdots | v_i) \otimes
(v_{i+1} | \cdots | v_n)$.  This involves no integrals, no
propagators, and no curve geometry.  Coassociativity
$(\Delta \otimes \id) \circ \Delta
= (\id \otimes \Delta) \circ \Delta$ is an identity of graded
vector spaces, valid at every genus.

(iii): The $\SCchtop$ operad has closed-color spaces
$C_\ast(\FM_k(\C))$ and open-color spaces $E_1(m)$.  At
genus $g \ge 1$, replacing $\C$ by $\Sigma_g$: the space
$\FM_k(\Sigma_g)$ is no longer contractible (it has
$H^1(\Sigma_g)^{\oplus k} \ne 0$), and the propagator acquires
quasi-periodic monodromy from the period matrix.  This is
the Arnold defect, and it contributes curvature.  The space
$\R$, however, does not change with genus: line operators
are still ordered on the boundary half-line $\R_{\ge 0}$,
independent of $\Sigma_g$.  Hence $E_1(m)$ is genus-independent.

(iv): For the affine lineage (where the Feigin--Frenkel
involution $k \mapsto -k - 2\hvee$ ensures
$\kappa(\cA) + \kappa(\cA^!) = 0$; Vol~I, Theorem~C
specialized to Kac--Moody and free-field algebras),
the closed-color curvatures
$\kappa(\cA) \cdot \omega_g$ and
$\kappa(\cA^!) \cdot \omega_g = -\kappa(\cA) \cdot \omega_g$
cancel when one passes to the full bulk theory controlled by
the Koszul pair $(\cA, \cA^!)$.  For $\mathcal{W}$-algebras
the sum $\kappa + \kappa^!$ is in general nonzero
(Vol~I, Theorem~D).
\end{proof}

\begin{remark}[Operadic explanation of the asymmetry]
\label{rem:asymmetry-operadic}
\index{genus tower!operadic explanation}
The asymmetry between colors at $g \ge 1$ is a direct
consequence of the operad structure.  The closed-color factor
$\FM_k(\Sigma_g)$ is a non-trivial fibration over
$\mathcal{M}_g$ with fibers that depend on the complex
structure; at $g = 0$, $\FM_k(\C)$ is a $K(\pi, 1)$ for the
pure braid group, and the bar differential is flat.  At
$g \ge 1$, the cohomology of $\FM_k(\Sigma_g)$ acquires new
classes from $H^1(\Sigma_g)$, and the bar differential squares
to $\kappa \cdot \omega_g \ne 0$.  The open-color factor
$E_1(m) \simeq \Sigma_m$ (a discrete set of orderings) is
independent of all continuous geometry and therefore of genus.
Curvature enters the two-color picture only through the
closed-color door.
\end{remark}

\begin{remark}[$\Theta_\cA$ and two colors]
\label{rem:theta-two-colors}
\index{shadow Postnikov tower!two-color decomposition}
\index{Maurer--Cartan element!two-color decomposition}
The universal Maurer--Cartan element
$\Theta_\cA \in \MC(\mathfrak{g}_\cA^{\mathrm{mod}})$
(Vol~I, Theorem~\ref*{V1-thm:mc2-bar-intrinsic}) encodes
\emph{both} colors simultaneously:
\begin{itemize}
\item \emph{Closed-color projection:} the shadow Postnikov
  tower $\kappa$, $C$, $Q$, \ldots are the finite-order
  projections of $\Theta_\cA$ onto the closed-color
  (chiral/modular) sector.  These are the invariants of the
  chiral bar complex $\barBch(\cA)$.
\item \emph{Open-color projection:} the $R$-matrix
  $r(z) = \Res^{\mathrm{coll}}_{0,2}(\Theta_\cA)$
  is the collision residue of $\Theta_\cA$ at arity~$2$
  in the open-color sector.  This is the invariant of the
  ordered bar complex $\barB^{\mathrm{ord}}(\cA)$.
\item \emph{Full MC equation:} $D \cdot \Theta_\cA
  + \tfrac{1}{2}[\Theta_\cA, \Theta_\cA] = 0$
  is a \emph{simultaneous} constraint on both colors.  The
  closed-color component gives $d_{\barB}^2 = 0$ (at
  genus~$0$) or $d_{\barB}^2 = \kappa \cdot \omega_g$ (at
  genus~$g \ge 1$).  The open-color component gives
  coassociativity of $\Delta$.  The mixed component gives
  the Yang--Baxter equation for $R(z)$.
\end{itemize}
The two-color Koszul duality theorem
(Theorem~\ref{thm:two-color-master}) unpacks the single MC
equation into its color components, making explicit the
information that $\Theta_\cA$ encodes simultaneously about
the chiral algebra and its line operators.
\end{remark}


%%%%%%%%%%%%%%%%%%%%%%%%%%%%%%%%%%%%%%%%%%%%%%%%%%%%%%%%%%%%%%%%%%%%%%%%%%%%%
% ELLIPTIC SPECTRAL DICHOTOMY AT GENUS 1
%%%%%%%%%%%%%%%%%%%%%%%%%%%%%%%%%%%%%%%%%%%%%%%%%%%%%%%%%%%%%%%%%%%%%%%%%%%%%

\subsection{Elliptic spectral dichotomy at genus~$1$}
\label{subsec:elliptic-spectral-dichotomy}
\index{elliptic spectral dichotomy}
\index{curvature!braiding coupling}
\index{genus 1!spectral braiding}

The genus-$1$ curvature $\dfib^{\,2} = \kappa(\cA)\cdot\omega_1$
(Theorem~\ref{thm:mc5-genus-one-bridge}) and the genus-$0$ spectral
$R$-matrix $R(z)$ (Theorem~\ref{thm:spectral_R_YBE}) originate from
different genera and appear to be independent structures.  At genus~$1$,
however, they meet: the spectral braiding must be defined on the
elliptic curve $E_\tau$, and the quasi-periodicity of the elliptic
propagator forces the two structures to interact.  The nature of this
interaction is controlled entirely by the pole structure of the
$\lambda$-bracket.

\begin{theorem}[Elliptic spectral dichotomy;
\ClaimStatusProvedHere]
\label{thm:elliptic-spectral-dichotomy}
Let $\cA$ be a logarithmic $\SCchtop$-algebra on $E_\tau \times \R$, and write the
$(-1)$-shifted $\lambda$-bracket on cohomology as
\[
\{a {}_\lambda b\}
\;=\;
c_0(a,b) + c_1(a,b)\,\lambda + c_2(a,b)\,\lambda^2 + \cdots\,,
\]
so that the genus-$0$ classical $r$-matrix
\textup{(}Proposition~\textup{\ref{prop:field-theory-r})} has the
Laurent expansion
\[
r(z)
\;=\;
\frac{c_0}{z} + \frac{c_1}{z^2} + \frac{c_2}{z^3} + \cdots
\qquad \text{\textup{(}formal, $z \to \infty$\textup{)}}.
\]
At genus~$1$, the classical $r$-matrix on $E_\tau$ is
\begin{equation}
\label{eq:elliptic-r-matrix}
r^{E_\tau}(z)
\;=\;
c_0 \cdot \zeta(z|\tau)
\;+\;
c_1 \cdot \wp(z|\tau)
\;+\;
\sum_{n \ge 2}
  \frac{(-1)^{n-1}}{(n-1)!}\,
  c_n \cdot \wp^{(n-2)}(z|\tau),
\end{equation}
obtained by replacing the rational functions $1/z^{n+1}$ in the
Laurent expansion of $r(z)$ by their elliptic counterparts on
$E_\tau$: the unique functions matching $1/z^{n+1} + O(z)$
near $z = 0$ and having prescribed periodicity.  Here
$\zeta(z|\tau)$, $\wp(z|\tau)$ are the Weierstrass zeta and
$\wp$-functions with periods $1, \tau$.

Then:
\begin{enumerate}[label=\textup{(\alph*)}]
\item \label{item:dichotomy-monodromy}
\textup{(Quasi-periodicity.)}
The $B$-cycle monodromy of $r^{E_\tau}$ is
\begin{equation}
\label{eq:r-monodromy}
r^{E_\tau}(z + \tau) - r^{E_\tau}(z)
\;=\;
2\eta_\tau \cdot c_0,
\end{equation}
where $\eta_\tau = \zeta(\tau/2|\tau)$ is the quasi-period.  The
$A$-cycle monodromy is $2\eta_1 \cdot c_0$.

\item \label{item:dichotomy-decoupling}
\textup{(Decoupling: Cartan type.)}
If $c_0 = 0$ (equivalently, if
$\{a {}_0 b\} = 0$ for all
$a, b \in H^\bullet(\cA, Q)$), then $r^{E_\tau}(z)$ is a
meromorphic function on $E_\tau$ \textup{(}doubly
periodic\textup{)}.
The spectral braiding is independent of the $B$-cycle monodromy
that sources the curvature $\kappa(\cA) \cdot \omega_1$.  The
genus-$1$ Swiss-cheese structure splits:
\[
\text{curved}\;\mathsf{SC}^{\mathrm{ch,top},(1)}
\;\simeq\;
(\text{curved closed color})
\;\times\;
(\text{flat open color with meromorphic braiding}).
\]

\item \label{item:dichotomy-entanglement}
\textup{(Entanglement: Yangian type.)}
If $c_0 \ne 0$, then $r^{E_\tau}(z)$ is quasi-periodic with
non-trivial $B$-cycle monodromy $2\eta_\tau \cdot c_0$.  The
spectral braiding is a section of a flat connection on the universal
cover of $E_\tau$ with holonomy~$c_0$, and the genus-$1$
Swiss-cheese structure does \emph{not} split.
The braiding and the curvature share a common geometric source:
the $B$-cycle monodromy of the Weierstrass zeta function.
\end{enumerate}
\end{theorem}

\begin{proof}
\textbf{Step~1} (Elliptic propagator in the exchange diagram).
The spectral $R$-matrix $R(z)$ at genus~$0$ is computed from
a single bulk-to-boundary propagator exchange
(Proposition~\ref{prop:field-theory-r}, Step~1), with
holomorphic propagator $K^{\C}(z) = dz/z$.  At genus~$1$,
the theory on $E_\tau \times \R_{\ge 0}$ uses the elliptic
propagator
$K^{E_\tau}(z) = \zeta(z|\tau)\,dz$
(Theorem~\ref{thm:mc5-genus-one-bridge}, Step~1), which has
the same local singularity
$\zeta(z|\tau) = 1/z + O(z)$ near $z = 0$
but acquires the quasi-periodicity
\[
\zeta(z + 1|\tau) = \zeta(z|\tau) + 2\eta_1,
\qquad
\zeta(z + \tau|\tau) = \zeta(z|\tau) + 2\eta_\tau.
\]

\textbf{Step~2} (Laplace formula at genus~$1$).
The $\lambda$-bracket on cohomology is a \emph{local} object:
it depends only on the residue of the propagator at the
collision divisor $z_1 = z_2$
(Section~\ref{sec:Ainfty-to-PVA}), which is the same at
genus~$0$ and genus~$1$
(Theorem~\ref{thm:mc5-genus-one-bridge}, Step~1).  Therefore
$\{a {}_\lambda b\}^{E_\tau} = \{a {}_\lambda b\}^{\C}$,
and the coefficients $c_n$ are genus-independent.

However, the classical $r$-matrix is a \emph{global} object:
it depends on the propagator at \emph{separated} points
$z_1, z_2$.  At genus~$1$, the Laplace formula of
Proposition~\ref{prop:field-theory-r} uses the elliptic
propagator $\zeta(z|\tau)$ in place of $1/z$.  The formal
substitution $1/z^{n+1} \leadsto f_{n+1}(z|\tau)$ is
determined by matching poles:
\begin{align*}
\frac{1}{z}
  &\;\leadsto\;
  \zeta(z|\tau)
  &&\text{(quasi-periodic)}, \\
\frac{1}{z^2}
  &\;\leadsto\;
  \wp(z|\tau) = -\zeta'(z|\tau)
  &&\text{(periodic)}, \\
\frac{1}{z^{n+1}}
  &\;\leadsto\;
  \frac{(-1)^{n}}{n!}\,\wp^{(n-1)}(z|\tau)
  &&\text{(periodic, $n \ge 2$)}.
\end{align*}
These are the unique functions with pole $1/z^{n+1} + O(z)$ at
$z = 0$ that are doubly periodic for $n \ge 1$, and
quasi-periodic with prescribed quasi-periods for $n = 0$.
Substituting into $r(z) = \sum_{n \ge 0} c_n / z^{n+1}$
gives equation~\eqref{eq:elliptic-r-matrix}.

\textbf{Step~3} (Quasi-periodicity).
The monodromy~\eqref{eq:r-monodromy} follows directly:
the only non-periodic function in the
expansion~\eqref{eq:elliptic-r-matrix} is
$c_0 \cdot \zeta(z|\tau)$, since all terms involving
$\wp^{(n-1)}$ for $n \ge 1$ are doubly periodic.
The quasi-periodicity
$\zeta(z + \tau|\tau) - \zeta(z|\tau) = 2\eta_\tau$
gives~\ref{item:dichotomy-monodromy}.

\textbf{Step~4} (Common geometric source).
The same $B$-cycle quasi-period $2\eta_\tau$ of the Weierstrass
zeta function produces two distinct algebraic structures:
\begin{itemize}
\item \emph{Curvature.}
In the proof of Theorem~\ref{thm:mc5-genus-one-bridge}
(Step~3), the quasi-periodicity of $\zeta$ breaks the Arnold
relation on $\FM_3(E_\tau)$, producing the defect
$E_2(\tau) \cdot (dz_1 - dz_2) \wedge (dz_2 - dz_3)$.
The Eisenstein series $E_2(\tau)$ is related to the
quasi-period by $\eta_1 = \pi^2 E_2(\tau) / 6$.
Contracting with the OPE data extracts the
\emph{highest-pole coefficient} of the $\lambda$-bracket
(the term $c_1 \cdot \lambda$ for affine algebras), yielding
the scalar curvature
$\dfib^{\,2} = \kappa(\cA) \cdot \omega_1$.

\item \emph{Braiding monodromy.}
In the elliptic $r$-matrix~\eqref{eq:elliptic-r-matrix},
the same quasi-periodicity produces the monodromy
$\delta r = 2\eta_\tau \cdot c_0$, which extracts the
\emph{lowest-pole coefficient} $c_0 = \{a {}_0 b\}$ of the
$\lambda$-bracket.
\end{itemize}
The curvature is a \emph{scalar} (the trace of the
highest-pole coefficient), invisible to the tensor structure
of the $r$-matrix.
The braiding monodromy is a \emph{tensor} (the full
lowest-pole coefficient), carrying the Lie-algebraic structure
constants.  When $c_0 = 0$, the braiding monodromy vanishes
and the only effect of the $B$-cycle is the scalar curvature:
decoupling~\ref{item:dichotomy-decoupling}.  When
$c_0 \ne 0$, both scalar and tensorial effects are sourced
by the same geometric defect:
entanglement~\ref{item:dichotomy-entanglement}.
\end{proof}

\begin{remark}[The Coisson bracket as discriminant]
\label{rem:Coisson-discriminant}
\index{Coisson bracket!curvature-braiding discriminant}
The coefficient $c_0 = \{a {}_0 b\}$ is the $\lambda = 0$
specialization of the $\lambda$-bracket, i.e.\ the Coisson bracket
of Theorem~\ref{thm:pva-coisson-bridge}.  The elliptic spectral
dichotomy is therefore:
\[
\text{Coisson bracket vanishes}
\;\;\Longleftrightarrow\;\;
\text{curvature-braiding decoupling at genus~$1$}.
\]
The Coisson structure is therefore the obstruction to the independence of holomorphic
and topological deformations at genus~$1$.
\end{remark}

\begin{remark}[Verification in the standard landscape]
\label{rem:elliptic-dichotomy-examples}
The dichotomy is verified in the standard examples:
\begin{enumerate}[label=(\roman*)]
\item \emph{Heisenberg $\mathcal{H}_k$:}
  $\{J {}_\lambda J\} = k\lambda$, so $c_0 = 0$.
  The elliptic $r$-matrix is
  $r^{E_\tau}(z) = k \cdot \wp(z|\tau)$,
  doubly periodic on $E_\tau$.  Curvature
  $\dfib^{\,2} = k \cdot \omega_1$
  (Theorem~\ref{thm:rosetta-curved}) and braiding decouple.

\item \emph{Affine Kac--Moody $\widehat{\fg}_k$:}
  $\{J^a {}_\lambda J^b\} = f^{ab}_c\, J^c +
  k\,\kappa^{ab}\,\lambda$,
  so $c_0^{ab} = f^{ab}_c\, J^c \ne 0$.
  The elliptic $r$-matrix
  $r^{E_\tau}(z) = \Omega \cdot \zeta(z|\tau) +
  k\kappa \cdot \wp(z|\tau)$
  has $B$-cycle monodromy $2\eta_\tau \cdot \Omega$
  (where $\Omega = \sum_a t^a \otimes t_a$ is the quadratic
  Casimir).  Curvature and braiding entangle.

\item \emph{Virasoro $\mathrm{Vir}_c$:}
  $\{L {}_\lambda L\} = (\partial + 2\lambda)\,L +
  \tfrac{c}{12}\,\lambda^3$,
  so $c_0 = \partial L \ne 0$.
  Curvature and braiding entangle.
\end{enumerate}
\end{remark}

\begin{remark}[Connection to the elliptic KZ equation]
\label{rem:elliptic-KZ}
For $\widehat{\fg}_k$, the elliptic $r$-matrix
$r^{E_\tau}(z) = \Omega \cdot \zeta(z|\tau) + \cdots$
is the classical limit of the elliptic Knizhnik--Zamolodchikov
connection of Bernard \cite{Ber88}:
\[
\nabla_i^{\mathrm{KZ}}
\;=\;
\partial_{z_i}
\;-\;
\frac{1}{k + h^\vee}
\sum_{j \ne i}
  \Omega_{ij} \cdot \wp(z_i - z_j|\tau).
\]
The quasi-periodicity of $\zeta$ in the $r$-matrix
corresponds to the non-trivial $B$-cycle monodromy of
the KZ connection.  The curvature-braiding entanglement
at genus~$1$ is, from this perspective, the statement
that the elliptic KZ connection on $E_\tau^n$ and the
modular differential equation
$\partial_\tau \Phi = \Delta_\tau \Phi$ (governing
genus-$1$ partition functions) share the same Casimir
data~$\Omega$.  For abelian algebras ($\Omega$ scalar),
the KZ connection is flat and trivializable;
for non-abelian algebras, the holonomy representation
of $\pi_1(E_\tau)$ is irreducible.
\end{remark}

\begin{remark}[Operadic splitting obstruction]
\label{rem:operadic-splitting}
In the language of the Swiss-cheese operad, the dichotomy
states: at genus~$1$, the curved $\SCchtop$-algebra structure
on $\barB(\cA)$ can be presented as a product of a curved
(closed-color) Beilinson--Drinfeld algebra and a flat (open-color)
$E_1$-algebra if and only if $c_0 = 0$.  When $c_0 \ne 0$,
the two colors are coupled by the non-trivial $B$-cycle
holonomy: the open-color braiding ``sees'' the closed-color
monodromy through the shared quasi-periodicity of the
propagator.  The obstruction to splitting is the element
$c_0 \in H^\bullet(\cA,Q)^{\otimes 2}$, the infinitesimal
generator of the holonomy representation.
\end{remark}
